%%%%%%%%%%%%%%%%%%%%%%%%%%%%%%%%%%%%%%%%%%%%%%%%%%%%%%%%%%%%%%%%%%%%%%%%%
%
% Plantilla para libro de texto de matemáticas.
%
% Esta plantilla ha sido desarrollada desde cero, pero utiliza algunas partes
% del código de la plantilla original utilizada en apuntesDGIIM
% (https://github.com/libreim/apuntesDGIIM), basada a su vez en las plantillas
% 'Short Sectioned Assignment' de Frits Wenneker (http://www.howtotex.com),
% 'Plantilla de Trabajo' de Mario Román y 'Plantilla básica de Latex en Español'
% de Andrés Herrera Poyatos (https://github.com/andreshp). También recoge
% ideas de la plantilla 'Multi-Purpose Large Font Title Page' de
% Frits Wenneker y Vel (vel@latextemplates.com).
%
% Licencia:
% CC BY-NC-SA 4.0 (https://creativecommons.org/licenses/by-nc-sa/4.0/)
%
%%%%%%%%%%%%%%%%%%%%%%%%%%%%%%%%%%%%%%%%%%%%%%%%%%%%%%%%%%%%%%%%%%%%%%%%%


% ---------------------------------------------------------------------------
% CONFIGURACIÓN BÁSICA DEL DOCUMENTO
% ---------------------------------------------------------------------------

%\documentclass[11pt, a4paper, twoside]{article} % Usar para imprimir
\documentclass[10pt, a4paper]{article}

\linespread{1.3}            % Espaciado entre líneas.
\setlength\parindent{0pt}   % No indentar el texto por defecto.
\setlength\parskip{7pt}

% ---------------------------------------------------------------------------
% CONFIGURACIÓN PERSONALIZADA
% ---------------------------------------------------------------------------

\usepackage{config}

% ---------------------------------------------------------------------------
% PAQUETES BÁSICOS
% ---------------------------------------------------------------------------

% IDIOMA
\usepackage[spanish, es-tabla, es-lcroman, es-noquoting, es-minimal]{babel}

% FUENTES
\usepackage{newpxtext, newpxmath}   % Fuente similar a Palatino
\usepackage{FiraSans}                 % Fuente sans serif
\usepackage[T1]{fontenc}
\usepackage[italic]{mathastext}     % Utiliza la fuente del documento
                                    % en los entornos matemáticos

% MÁRGENES
\usepackage[margin=2.5cm, top=3cm]{geometry}

%  OTROS
\usepackage{hyperref}   % Enlaces
\usepackage{graphicx}   % Permite incluir gráficos en el documento

% ---------------------------------------------------------------------------
% FUENTE DE MATEMÁTICAS ADAPTABLE EN SANS-SERIF
% ---------------------------------------------------------------------------

\usepackage{sansmath}

\newif\IfInSansMode
\let\oldsf\sffamily
\renewcommand*{\sffamily}{\oldsf\sansmath\InSansModetrue}

% ---------------------------------------------------------------------------
% DISEÑO DE PÁGINA
% ---------------------------------------------------------------------------

\usepackage{pagecolor}
\usepackage{afterpage}

% ---------------------------------------------------------------------------
% CABECERA Y PIE DE PÁGINA
% ---------------------------------------------------------------------------

\usepackage{fancyhdr}   % Paquete para cabeceras y pies de página

% Indica que las páginas usarán la configuración de fancyhdr
\pagestyle{fancy}
\fancyhf{}

% Representa la sección de la cabecera
\renewcommand{\sectionmark}[1]{%
\markboth{#1}{}}

% Representa la subsección de la cabecera
\renewcommand{\subsectionmark}[1]{%
\markright{#1}{}}

% Parte derecha de la cabecera
\fancyhead[LE,RO]{\sffamily\textsl{\rightmark} \hspace{1em}  \textcolor{500}{\rule[-0.4ex]{0.2ex}{1.2em}} \hspace{1em} \thepage}

% Parte izquierda de la cabecera
\fancyhead[RE,LO]{\sffamily{\leftmark}}

% Elimina la línea de la cabecera
\renewcommand{\headrulewidth}{0pt}

% Controla la altura de la cabecera para que no haya errores
\setlength{\headheight}{14pt}

% ---------------------------------------------------------------------------
% TÍTULOS DE PARTES Y SECCIONES
% ---------------------------------------------------------------------------

\usepackage{titlesec}

% Estilo de los títulos de las partes
\titleformat{\part}[hang]{\Huge\bfseries\sffamily}{\thepart\hspace{20pt}\textcolor{500}{|}\hspace{20pt}}{0pt}{\Huge\bfseries}
\titlespacing*{\part}{0cm}{-2em}{2em}[0pt]

% Reiniciamos el contador de secciones entre partes (opcional)
\makeatletter
\@addtoreset{section}{part}
\makeatother

% Estilo de los títulos de las secciones, subsecciones y subsubsecciones
\titleformat{\section}
  {\Large\bfseries\sffamily}{\thesection}{1em}{}

\titleformat{\subsection}
  {\Large\sffamily}{\thesubsection}{1em}{}[\vspace{.5em}]

\titleformat{\subsubsection}
  {\sffamily}{\thesubsubsection}{1em}{}

% ---------------------------------------------------------------------------
% ENTORNOS PERSONALIZADOS
% ---------------------------------------------------------------------------

\usepackage{mdframed}

%% MARCOS DE LOS ESTILOS

% Configuración general de mdframe, los estilos de los teoremas, etc
\mdfsetup{
  skipabove=1em,
  skipbelow=1em,
  innertopmargin=1em,
  innerbottommargin=1em,
  splittopskip=2\topsep,
}

% Definimos los marcos de los estilos

\mdfdefinestyle{nth-frame}{
	linewidth=2pt, %
	linecolor= 500, %
	topline=false, %
	bottomline=false, %
	rightline=false,%
	leftmargin=0em, %
	innerleftmargin=1em, %
  innerrightmargin=1em,
	rightmargin=0em, %
}%

\mdfdefinestyle{nprop-frame}{
	linewidth=2pt, %
	linecolor= 300, %
	topline=false, %
	bottomline=false, %
	rightline=false,%
	leftmargin=0pt, %
	innerleftmargin=1em, %
	innerrightmargin=1em,
	rightmargin=0pt, %
}%

\mdfdefinestyle{ndef-frame}{
	linewidth=2pt, %
	linecolor= 500, %
	backgroundcolor= 50,
	topline=false, %
	bottomline=false, %
	rightline=false,%
	leftmargin=0pt, %
	innerleftmargin=1em, %
	innerrightmargin=1em,
	rightmargin=0pt, %
}%

\mdfdefinestyle{ejer-frame}{
	linewidth=2pt, %
	linecolor= 300, %
	backgroundcolor= 50,
	topline=false, %
	bottomline=false, %
	rightline=false,%
	leftmargin=0pt, %
	innerleftmargin=1em, %
	innerrightmargin=1em,
	rightmargin=0pt, %
}%

\mdfdefinestyle{ejemplo-frame}{
	linewidth=0pt, %
	linecolor= 300, %
	leftline=false, %
	rightline=false, %
	leftmargin=0pt, %
	innerleftmargin=1.3em, %
	innerrightmargin=1em,
	rightmargin=0pt, %
	innertopmargin=0em,%
	innerbottommargin=0em, %
	splittopskip=\topskip, %
}%

% Asignamos los marcos a los estilos
\surroundwithmdframed[style=nth-frame]{nth}
\surroundwithmdframed[style=nprop-frame]{nprop}
\surroundwithmdframed[style=nprop-frame]{ncor}
\surroundwithmdframed[style=ndef-frame]{ndef}
\surroundwithmdframed[style=ejer-frame]{ejer}
\surroundwithmdframed[style=ejemplo-frame]{ejemplo}
\surroundwithmdframed[style=ejemplo-frame]{sol}

%%%%%%%%%%%%%%%%%%%%%%%%%%%%%%%%%%%%%%%%%%%%%%%%%%%%%%%%%%%%%%%%%%%%%%%%%%%%%
% ---------------------------------------------------------------------------
% COMIENZO DEL DOCUMENTO
% ---------------------------------------------------------------------------
%%%%%%%%%%%%%%%%%%%%%%%%%%%%%%%%%%%%%%%%%%%%%%%%%%%%%%%%%%%%%%%%%%%%%%%%%%%%%

\begin{document}

% ---------------------------------------------------------------------------
% PORTADA EXTERIOR
% ---------------------------------------------------------------------------

\newpagecolor{500}\afterpage{\restorepagecolor} % Color de la página
\begin{titlepage}

  % Título del documento
	\parbox[t]{\textwidth}{
			\raggedright % Texto alineado a la izquierda
			\fontsize{50pt}{50pt}\selectfont\sffamily\color{white}{
			  \textbf{\asignatura}
      }
	}

	\vfill

	%% Autor e información del documento
	\parbox[t]{\textwidth}{
		\raggedright % Texto alineado a la izquierda
		\sffamily\large\color{white}
		{\Large \autor }\\[4pt]
		\grado\\
		\universidad\\[4pt]
		\texttt{\enlaceweb}
	}

\end{titlepage}

% ---------------------------------------------------------------------------
% PÁGINA DE LICENCIA
% ---------------------------------------------------------------------------

\thispagestyle{empty}
\null
\vfill

%% Información sobre la licencia
\parbox[t]{\textwidth}{
  \includegraphics{by-nc-sa.pdf}\\[4pt]
  \raggedright % Texto alineado a la izquierda
  \sffamily\large
  {\Large Este libro se distribuye bajo una licencia CC BY-NC-SA 4.0.}\\[4pt]
  Eres libre de distribuir y adaptar el material siempre que reconozcas a los\\
  autores originales del documento, no lo utilices para fines comerciales\\
  y lo distribuyas bajo la misma licencia.\\[4pt]
  \texttt{creativecommons.org/licenses/by-nc-sa/4.0/}
}

% ---------------------------------------------------------------------------
% PORTADA INTERIOR
% ---------------------------------------------------------------------------

\begin{titlepage}

  % Título del documento
	\parbox[t]{\textwidth}{
			\raggedright % Texto alineado a la izquierda
			\fontsize{50pt}{50pt}\selectfont\sffamily\color{500}{
			  \textbf{\asignatura}
      }
	}

	\vfill

	%% Autor e información del documento
	\parbox[t]{\textwidth}{
		\raggedright % Texto alineado a la izquierda
		\sffamily\large
		{\Large \autor}\\[4pt]
		\grado\\
		\universidad\\[4pt]
		\texttt{\enlaceweb}
	}

\end{titlepage}

% ---------------------------------------------------------------------------
% ÍNDICE
% ---------------------------------------------------------------------------

\thispagestyle{empty}
\tableofcontents
\newpage

% ---------------------------------------------------------------------------
% CONTENIDO
% ---------------------------------------------------------------------------

\part{Teoría}
% Apuntes originales por Santiago de Diego, Braulio Valdivielso y Francisco Luque
% https://github.com/santidediego/Probabilidad-I-DGIIM

\chapter{Conjuntos y funciones $\sigma$-aditivas}

\section{Introducción: Espacio medible}

\subsection*{Concepto de $\sigma$-álgebra}
Definamos primero la estructura de álgebra. Dado un espacio $\Omega$, una clase de subconjuntos de $\Omega$, $\mathcal{Q} \subset P(\Omega)$ tiene estructura de álgebra si y solo si:
\begin{enumerate}
	\item $\Omega \in \mathcal{Q}$
	\item $\forall A \in \mathcal{Q} \implies A^c \in \mathcal{Q}$
	\item $\forall A,B \in \mathcal{Q} \implies A \cup B \in \mathcal{Q}$
\end{enumerate}


Una vez vista esta definición, sea un conjunto $\Omega$, una
familia de subconjuntos $\mathcal{A}$ es $\sigma$-álgebra sobre $\Omega$ si: 
\begin{enumerate}
\item $\Omega \in \mathcal{A}$
\item Si $A \in \mathcal{A} \Rightarrow \overline{A} \in \mathcal{A}$
\item Si $A_1, A_2, A_3\ldots \in \mathcal{A} \Rightarrow \displaystyle \bigcup_{n \in \mathbb{N}}^{\infty} A_n \in \mathcal{A}$
\end{enumerate}
Una $\sigma$-álgebra es una clase de conjuntos cerrada para el complementario y unión numerable. Existen una serie de propiedades inmediatas derivadas:
\begin{itemize}
  \item $\emptyset \in \mathcal{A}$
  \item $\mathcal{A}$ es cerrada para la  intersección numerable.
\end{itemize}

\subsection*{Espacio medible}

Al par ($\Omega$, $\mathcal{A}$) se le llama espacio medible, y a los conjuntos que pertenecen a $\mathcal{A}$ se les denomina conjuntos medibles.

\section{Límites en sucesiones de conjuntos}

\subsection{Límites en sucesiones monótonas}

Si tomamos la relación de inclusión entre conjuntos ($\subseteq$) como una relación de orden, podemos hablar sin ningún tipo de problema de sucesiones monótonas. En este sentido, una sucesión   $\{A_n\}_{n \in \mathbb{N}}$  creciente de conjuntos cumpliría que $A_i \subseteq A_{i+1}, \forall i \in \mathbb{N}$. De forma análoga para sucesión decreciente de conjuntos.\\

El límite de una sucesión creciente de conjuntos $\{A_n\}_{n \in \mathbb{N} }$ se puede definir
$$\lim_{n\to\infty} A_n = \bigcup_{k=1}^{\infty} A_k$$

Esta definición aprovecha la relación de orden entre los conjuntos de la sucesión para afirmar que \textit{si un elemento está en el último conjunto de la sucesión, entonces está en todos los anteriores, y por tanto en \textbf{todos}}, por ello se puede utilizar una intersección numerable (que está perfectamente definida) para formalizar el concepto.\\

Con una intuición análoga se define el límite de una sucesión decreciente de conjuntos. Si $\{A_n\}_{n\in\mathbb{N}}$ es una sucesión decreciente de conjuntos, entonces:

$$ \lim_{n\to\infty} A_n = \bigcap_{k=1}^{\infty} A_k $$

\subsection{Límites superiores e inferiores}

No solo se puede hablar de límites en sucesiones monótonas. 
Para definir los límites en sucesiones arbitrarias de conjuntos tenemos que recurrir a los conceptos de límite inferior y superior. 
La intuición de estos límites superior e inferior pasan por el concepto de las colas de la sucesión.\\

Dada una sucesión de conjuntos $\{A_n\}_{n\in\mathbb{N} }$ podemos considerar el conjunto

$$ B_n = \bigcap_{k=n}^{\infty} A_k $$

y este conjunto contiene aquellos elementos que están en \textbf{todos} los $A_k$ para $k \geq n$. Es fácil probar que la sucesión $\{B_n\}_{n\in\mathbb{N}}$ es una sucesión creciente de conjuntos, y por tanto se puede obtener el límite $\lim_{n\to\infty} B_n$. Informalmente, ese límite es un conjunto que contiene a todos los elementos de $A_n$ que están en todos los conjuntos $A_k$ a partir de cierto $n\in\mathbb{N}$. Definimos el límite inferior de $A_n$ como

$$ \liminf A_n = \lim_{n\to\infty} B_n = \bigcup_{n\in\mathbb{N}} \bigcap_{k=n}^{\infty} A_n $$

Análogamente, podríamos definir la sucesión de conjuntos $\{C_n\}_{n\in\mathbb{N}}$ como

$$ C_n = \bigcup_{k=n}^{\infty} A_n $$

Cada $C_n$ contiene todos los elementos que están presentes en algún $A_k$ para $k \geq n$. Es fácil también ver que la sucesión $\{C_n\}_{n\in\mathbb{N}}$ es una sucesión decreciente de conjuntos, y por tanto su límite también está bien definido. Se puede definir entonces el límite superior de $\{A_n\}_{n\in\mathbb{N}}$ como

$$ \limsup A_n = \lim_{n\to\infty} C_n = \bigcap_{n\in\mathbb{N}} \bigcup_{k=n}^{\infty} A_n $$

Informalmente se puede pensar en este límite superior de $A_n$ como el conjunto de los elementos que están en infinitos conjuntos de la sucesión.\\

A partir de estas definiciones, es fácil comprobar que

$$ \liminf A_n \subseteq \limsup A_n $$

Basta ver que un $\omega$ cualquiera de $\liminf A_n$ está en $\limsup A_n$. En efecto,
\[
\omega \in \liminf A_n \implies \omega \in \cup_{k=1}^\infty \cap_{n=k}^\infty A_n \implies \omega \in \cap_{n=k}^\infty A_n
\]
para algún $k$ finito, luego $\omega$ pertenece a una infinidad de $A_n$ y así:
\[
\omega \in \limsup A_n 
\]
Por lo que tenemos el resultado.

\subsection{Límite de una sucesión de conjuntos}

Diremos que una sucesión de conjuntos tiene límite si su límite inferior y superior coinciden, y el límite tendrá dicho  valor. Es decir,

$$ \lim A_n = \liminf A_n = \limsup A_n $$

en caso de que límite superior e inferior coincidan.

\section{Funciones sobre conjuntos}

Una vez definidos los conceptos sobre conjuntos con los que vamos a trabajar, pasamos a definir las funciones sobre conjuntos. Vamos entonces a definir lo que es una función de conjunto. Sean los espacios medibles ($\Omega$, $\mathcal{A}$) y($\Omega'$, $\mathcal{A}'$), definimos la función:
$$ X: \mathcal{A} \to \mathcal{A}'$$
$$ A \longrightarrow X(A)$$

A raíz de esta definición podemos definir también lo que se conoce como función inversa. Dada una función $X$, la función inversa de $X$, $X^{-1}$, asigna a cada conjunto $A' \in \mathcal{A}'$ el conjunto $A \in \mathcal{A}$ tal que $X(A) = A'$. La propiedad básica que cumplen las funciones inversas es que preservan las operaciones e inclusiones de conjuntos.

\section{Concepto de $\sigma$-aditividad}
Sea un conjunto $\Omega$ y una $\sigma$-álgebra $\mathcal{A}$ sobre $\Omega$. Definimos la función de conjunto:
$$\varphi : \mathcal{A} \rightarrow \mathbb{R}$$
\begin{center}
$\varphi (A)$ es único\\
\end{center}
Se dice que $\varphi$ es aditiva si
\[
\forall \{A_i\}_{i=1,...,n} \in A \text{ con } A_i \cap A_j = \emptyset, \ i \ne j , \varphi( \cup_{i=1}^n A_i ) = \sum_{i=1}^n \varphi(A_i)
\]

Se dice que $\varphi$ es $\sigma$-aditiva si
\[
\forall \{A_i\}_{n\in N}, \ A_n \in A \text{ con } A_i \cap A_j = \emptyset, \ i \ne j , \varphi( \cup_{i=1}^\infty A_n ) = \sum_{i=1}^\infty \varphi(A_n)
\]

\subsubsection{Función subaditiva}
Sea $\varphi$ definida como antes. Se dice que $\varphi$ es \textbf{subaditiva} si $\varphi(A \cup B) \leq \varphi(A) + \varphi(B)$

\begin{lemma}
Si $\exists B\, : \, \varphi(B)<\infty \Rightarrow \varphi(\emptyset)=0$
\end{lemma}
\begin{proof}
  $\varphi(B) =  \varphi(B+\emptyset) = \varphi(B) + \varphi(\emptyset)\implies \varphi(\emptyset) = 0$. Notese que $\varphi(B)<\infty$ es obligado porque si no $\varphi(\emptyset)$ podría tomar un valor $a \neq 0$ y se mantendría que  $\varphi(B) + \varphi(\emptyset) = \infty + a = \infty$.
\end{proof}


\begin{theorem}
Si $\varphi$ es finita, $\sigma$-aditiva y $\displaystyle|\sum\varphi(A_i)|<\infty\Rightarrow\displaystyle\sum(\vert \varphi(A_i)\vert)<\infty$
\end{theorem}
\begin{proof}
Tomemos las sucesiones: \\
Si $\varphi(A_n)\geq 0 \Longrightarrow A_n^{+}=A_n$ y $A_n^{-}=\emptyset$\\
Si $\varphi(A_n)<0 \Longrightarrow A_n^{+}=\emptyset$ y $A_n^{-}=A_n$\\
Entonces $\varphi(\displaystyle\sum A_n^{+})=\displaystyle\sum\varphi(A_n^{+})$ y $\varphi(\displaystyle\sum A_n^{-})=\displaystyle\sum\varphi(A_n^{-})$, ambas finitas. Sumando ambas cantidades obtenemos que $\displaystyle\sum(\vert \varphi(A_i)\vert)<\infty$


\end{proof}

\begin{theorem}
  \ 
  
Si $\varphi$ es $\sigma$-aditiva, $\varphi \geq 0$, entonces: $\varphi$ es no decreciente y subaditiva
\end{theorem}

\begin{proof}
  \
   
Veamos que es no decreciente $(A \subseteq B \Rightarrow \varphi (A) \leq \varphi (B))$. Podemos escribir $B = (B \cap A) + (B \cap \overline{A}) = A + (B \cap \overline{A})$. Entonces $\varphi (A) \leq \varphi (A) + \varphi (B \cap \overline{A}) = \varphi (B)$.

Veamos ahora que es subaditiva $(\varphi(\bigcup_{j} A_j) \le \sum_j \varphi(A_j))$.Veamos que 
$\varphi(\bigcup_{i} A_i) = \varphi(A_1 + A_1^c A_2 + ... + A_1^cA_2^c...A_n+ ...) =^1 \varphi(A_1) + \varphi(A_1^c A_2) + ... + \varphi(A_1^cA_2^c...A_n) + ... \leq  \sum_i \varphi(A_i)$ donde en 1 usamos que son disjuntos para aplicar la subaditividad.

\end{proof}

\section{Continuidad en funciones sobre conjuntos}

Una vez introducido el concepto de límite para una sucesión de conjuntos, vamos a tratar de definir la continuidad para funciones de conjunto. Tendremos tres tipos de continuidad, cada uno relacionado con un tipo de sucesiones de conjuntos de las que hemos definido anteriormente. En esta sección trabajaremos con una función $\varphi : \Omega \to \Omega'$\\

Diremos que $\varphi$ es continua por abajo si cumple que, dada una sucesión creciente de elementos $A_n \uparrow A$, se tiene que

$$ \lim \varphi (A_n) = \varphi (A) $$

Por otra parte, diremos que $\varphi$ es continua por arriba si $\exists A_{n_0} : \varphi (A_{n_0}) < \infty$ y cumple que dada una sucesión decreciente de elementos $A_n \downarrow A$, se tiene que

$$ \lim \varphi (A_n) = \varphi (A)$$

Por último, diremos que una función es continua si lo es por arriba y por abajo.

\begin{theorem}
Teorema de continuidad para funciones sobre conjuntos\\

Sea $\varphi$ una función $\sigma$-aditiva. Entonces, $\varphi$ es aditiva y continua. Inversamente, si $\varphi$ es aditiva y, o bien continua por abajo, o finita y continua en $\emptyset$, entonces $\varphi$ es $\sigma$-aditiva.
\end{theorem}

\begin{proof}
Por un lado, sea $\varphi$ una función $\sigma$-aditiva. Entonces es trivialmente aditiva. Ahora, veamos que es continua por abajo y por arriba. Sea $A_n \uparrow A$, entonces:
$$ A = \lim A_n = \bigcup A_n = A_1 + (A_2 - A_1) + (A_3 - A_2) +... $$

Unión de conjuntos disjuntos. Por tanto:
$$ \varphi (A) = \varphi (\lim A_n) = \lim_{n \to \infty } \{ \varphi (A_1) + \varphi (A_2 - A_1) + ... + \varphi (A_n - A_{n-1}) \} = \lim \varphi (A_n) $$

Veamos la continuidad por arriba. Sea $A_n \downarrow A$, tomamos $A_{n_0}$ tal que $\varphi (A_{n_0})$ es finito. Entonces $A_{n_0} - A_n \uparrow A_{n_0} - A$, y por el apartado anterior tenemos la convergencia desde abajo, por tanto:
$$ \varphi (A_{n_0}) - \varphi (A) = \varphi (\lim (A_{n_0} - A_n)) = \lim \varphi (A_{n_0} - A_n) = \varphi (A_{n_0}) - \lim \varphi (A_n) $$

De donde se deduce que $ \varphi (A) = \lim \varphi (A_n) $.\\

Inversamente, sea $ \varphi$ una función aditiva. Si $\varphi$ es continua por abajo, tenemos
$$ \varphi \left( \sum_{k=1}^{\infty} A_k \right) = \varphi \left( \lim \sum_{k=1}^n A_k \right) = \lim \varphi \left( \sum_{k=1}^n A_k \right) = \lim \sum_{k=1}^n \varphi \left( A_k \right) = \sum_{k=1}^{\infty} \varphi (A_k) $$

Y por tanto es $\sigma$-aditiva. Si es finita y continua en $\emptyset$, entonces se obtiene la $\sigma$-aditividad de:
$$ \varphi \left( \sum_{k=1}^{\infty} A_k \right) = \varphi \left( \sum_{k=1}^{n} A_k \right) + \varphi \left( \sum_{k=n+1}^{\infty} A_k \right) = \sum_{k=1}^{n} \varphi (A_k) + \varphi \left( \sum_{k=n+1}^{\infty} A_k \right) $$

Y tomando límites tenemos que
$$ \varphi \left( \sum_{k=n+1}^{\infty} A_k \right) \to \varphi (\emptyset) = 0 $$

\end{proof}

Una vez demostrado este teorema, vamos a ver un teorema que nos relaciona las propiedades del supremo e ínfimo de una función $\sigma$-aditiva con los conjuntos sobre los que está dicha función definida:

\begin{theorem}
Teorema del supremo e ínfimo\\

Sea $\varphi$ una función $\sigma$-aditiva sobre una $\sigma$-álgebra $\mathcal{A}$. Entonces, existen $C,D \in \mathcal{A}$ tales que $\varphi (C) = \sup \varphi$ y $\varphi (D) = \inf \varphi$
\end{theorem}

\begin{proof}
Probaremos la existencia del conjunto $C$. La del conjunto $D$ es análoga. Si $\varphi(A) = \infty$ para algún $A \in \mathcal{A}$, entonces podemos establecer $A = C$ y la demostración del teorema es trivial. Entonces, supongamos que $\varphi < \infty$ y dado que el valor $-\infty$ está excluido, $\varphi$ es finita.\\

Entonces, existe una sucesión $\{A_n\} \subset \mathcal{A}$ tal que $\varphi(A_n) \to \sup \varphi$. Sea $A = \cup A_n $ y para cada $n$, consideramos la partición de $A$ en $2^n$ conjuntos $A_{nm}$ de la forma $\displaystyle \cap_{k=1}^n A'_k$, donde $A'_k = A_k$ o $A - A_k$. Para $n < n'$, cada conjunto $A_{nm}$ es una suma finita de conjuntos $A_{n'm'}$. Sea ahora $B_n$ la suma de los conjuntos $A_{nm}$ para los cuales $\varphi$ es no negativa. Si no hay ninguno, entonces $B_n = \emptyset$. Entonces, tenemos por un lado que $A_n$ es la suma de algunos de los $A_{nm}$, y por otro lado, para $n' > n$, cada conjunto $A_{n'm'}$ está dentro de $B_n$ o son disjuntos. Entonces tenemos
$$ \varphi(A_n) \leq \varphi(B_n) \leq \varphi (B_n \cup B_{n+1} \cup ... \cup B_{n'}) $$
Tomando $n' \to +\infty$, se sigue de la continuidad por debajo que
$$ \varphi(A_n) \leq \varphi (B_n) \leq \varphi \left( \cup_{k=n}^{+\infty} B_k\right) $$
Haciendo tender ahora $n \to +\infty$ y tomando $\displaystyle C = \lim \cup_{k = n}^{+\infty} B_k$ se sigue, por la continuidad por arriba, que $\sup \varphi \leq \varphi(C)$. Pero $\varphi(C) \leq \sup \varphi$, y por tanto se tiene que $\varphi(C) = \sup \varphi$, como queríamos.

\end{proof}

\chapter{Medidas y probabilidades}

\section{Concepto de medida, media exterior y probabilidad}
Una función de conjuntos $\mu^\circ$ es una \textbf{medida} si verifica:
\begin{itemize}
\item Es $\sigma$-aditiva, es decir,  $\mu^\circ (\cup A_j) = \displaystyle \sum \mu^\circ (A_j)$
\item Es no decreciente $A\subset B \, \mu^\circ(A)\leq \mu^\circ(B)$
\item $\mu^\circ(\emptyset)=0$
\end{itemize}

A la tupla ($\Omega$, $\mathcal{A}$, $\mu_{\mathcal{A}}$) se le denomina espacio de medida.\\

Una \textbf{medida exterior} es una función de conjuntos que cumple:
\begin{itemize}
	\item Es positiva
	\item Es $\sigma$-subaditiva, es decir: $\forall \{A_n\}_{n\in N} \subset P(\Omega),$ se tiene:
	\[
	 \mu^e(\bigcup_{n=1}^\infty A_N) \leq \sum_{n=1}^\infty \mu^e(A_n)
	\]
\end{itemize} 

Para que una medida exterior fuera una medida tendría que ser $\sigma$-aditiva. Es decir, la falla la primera condición. Sí que es positiva ya que $\mu^\circ(\emptyset)=0$ y es creciente. Esta medida exterior se aplica a cualquier conjunto. Va a haber unos subconjuntos en los que la función se comporte como si fuera aditiva.\\

Una vez definido el concepto de medida, vamos a dar ahora el de probabilidad. Una probabilidad $\mathcal{P}$ sobre un espacio medible ($\Omega$, $\mathcal{A}$) es una medida que además cumple que $\mathcal{P}(\Omega)=1$. Por tanto, tiene las siguientes propiedades:\\

\begin{itemize}
	\item $\mathcal{P}(\emptyset)=0$
	\item $\forall A \in \mathcal{A}, 0 \leq \mathcal{P}(A) \leq 1$
	\item Es una función $\sigma$-aditiva
\end{itemize}

De forma análoga al concepto de espacio de medida, podemos definir ahora el de espacio probabilístico. Un espacio probabilístico es una tupla formada por un conjunto $\Omega$, una $\sigma$-álgebra sobre ese conjunto, $\mathcal{A}$, y una función de probabilidad $\mathcal{P}$.

\begin{theorem}
Teorema de sucesión
\begin{enumerate}
\item Si $A_n \uparrow A \Rightarrow \mathcal{P} (A_n) \uparrow \mathcal{P}(A)$
\item $\mathcal{P}(\lim \inf A_n) \leq \lim \inf \mathcal{P}(A_n)$
\item Si $A_n \downarrow A \Rightarrow \mathcal{P} (A_n) \downarrow \mathcal{P}(A)$
\item $\mathcal{P}(\lim \sup A_n) \geq \lim \sup \mathcal{P}(A_n)$
\item Si $A_n \rightarrow A \Rightarrow \mathcal{P} (A_n) \rightarrow \mathcal{P}(A)$
\end{enumerate}
\end{theorem}

\begin{proof}
\begin{enumerate} 
  \item Trivial por continuidad de la probabilidad.
  \item $\lim \inf (A_n) = (\cup_{n=1}^\infty\cap_{k=n}^\infty A_k)$.Sea $B_n := \cap_{k\geq n}^\infty A_k$. Entonces sabemos que $B_n \uparrow$ y que, por lo tanto $\lim B_n = \cup_n^\infty B_n = lim \inf (A_n)$. 
  
  Por último sabemos que $\mathcal{P}(A_n) \geq \mathcal{P}(B_n)$ dado que una probabilidad cumple ser  no decreciente.
  
  De todo esto podemos considerar que $\lim \inf \mathcal{P}(A_n) \geq \lim \mathcal{P}(B_n) = \mathcal{P}(\lim B_n) = \mathcal{P}(\lim \inf A_n)$
  
  \item trivial por continuidad de la probabilidad.
  
  \item Análogo 2.
  
  \item Se cumple por 2  y 4.
\end{enumerate}
\end{proof}

\section{Teorema de extensión de Carathéodory}

Una vez vistas las definiciones anteriores de medida y medida exterior, vamos a ver un teorema fundamental, que nos servirá para justificar la forma en la que haremos el cálculo de probabilidades a posteriori. Veamos dos lemas antes que nos serán necesarios para la demostración del teorema principal.\\

\textbf{Definición:} Si $\varphi, \varphi'$ son dos funciones de conjuntos sobre $\mathcal{A}$ y $\mathcal{A}'$ respectivamente, y toman los mismos valores en un subconjunto común $\mathcal{A}''$, se dice que coinciden en este subconjunto.

\textbf{Definición:} Si $\mathcal{A} \subset \mathcal{A}'$ y $\varphi, \varphi'$ coinciden en $\mathcal{A}$, decimos que $\varphi$ es una restricción de $\varphi'$, o que $\varphi'$ es una \textbf{extensión} de $\varphi$

\textbf{Definición:} Un conjunto $A\in S(\Omega)$ es $\mu^\circ$-medible si se cumple que $\mu^\circ(D) \geq \mu^\circ(AD)+\mu^\circ(A^cD), \forall D \in S(\Omega)$\\

\textbf{Definición:} Una extensión exterior $\mu^\circ$ de una medida $\mu$ se define como:
$$\mu^\circ (A) = \inf \sum_j \mu(A_j), \quad A \subset \cup A_j,\quad recubrimiento $$

\begin{theorem}
La extensión exterior $\mu^\circ$ de una medida $\mu$ sobre un álgebra es una extensión de $\mu$ a una medida exterior
\end{theorem}

\begin{proof}
Primero, veamos que es una extensión. Esto es trivial, dado que $\mu^\circ(A) = \inf \displaystyle \sum \mu(A) = \mu(A)$. Ahora, veamos que es una medida exterior.\\

Veamos que es subaditiva, es decir, $\mu^\circ(\displaystyle \cup_j A_j) \leq \sum_j \mu^\circ(A_j)$. Sea entonces, para cada $A_j$, un recubrimiento $ \displaystyle \cup_k A_{kj} \supset A_j$. Entonces:
$$ \sum_k \mu (A_{kj}) \leq \frac{\varepsilon}{2^j} + \mu^\circ (A_j)$$
$$ \mu^\circ (\cup_j A_j) \leq \sum_{j,k} \mu (A_{kj}) \leq \sum_j \Big( \mu^\circ(A_j) + \frac{\varepsilon}{2^j}\Big) \leq \sum_j (\mu^\circ(A_j)) + \varepsilon $$
Claramente es no decreciente, ya que si $A \subseteq B$, todo recubrimiento de $B$ lo será también de $A$, y por tanto, $\mu^\circ(A) \leq \mu^\circ(B)$\\

Por último, $\mu^\circ(\emptyset) = \mu(\emptyset) = 0$


\end{proof}

\begin{theorem}
Sea $\mu^\circ$ una medida exterior sobre la $\sigma$-álgebra $\mathcal{A}$, y sea $\mathcal{A}^\circ$ la clase de conjuntos $\mu^\circ$-medibles:
\begin{itemize}
\item $\mathcal{A}^\circ$ es un $\sigma$-campo
\item $\mu^\circ$ en $\mathcal{A}^\circ$ es una medida
\end{itemize}
\end{theorem}

\begin{proof}
Primero, veamos que $\mathcal{A}^\circ$ es un álgebra. Sea $A \in \mathcal{A}^\circ$, entonces $A^c \in \mathcal{A}^\circ$ dado que la definición de $\mu^\circ$-medibilidad es simétrica. Sean entonces $A,B \in \mathcal{A}^\circ$. Entonces:
$$ \mu^\circ(D) = \mu^\circ(AD) + \mu^\circ(A^cD) = $$
$$ = \mu^\circ(ABD) + \mu^\circ(AB^cD) + \mu^\circ(A^cBD) + \mu^\circ(A^cB^cD) \geq$$
$$ \geq \mu^\circ(ABD) + \mu^\circ (AB^cD \cup A^cBD \cup A^cB^cD) =$$
$$ = \mu^\circ(ABD) + \mu^\circ(AB)^cD$$
Dado entonces que $\mathcal{A}^\circ$ es cerrada bajo la operación complementario e intersección finita, lo es para la unión finita, así que nos encontramos ante un álgebra. Veamos ahora que $\mu^\circ$ es aditiva en $\mathcal{A}^\circ$. Sean $A,B \in \mathcal{A}^\circ$ disjuntos. Entonces:
$$\mu^\circ(A+B) = \mu^\circ((A+B)A) + \mu^\circ((A+B)A^c) = \mu^\circ(A) + \mu^\circ(B)$$
Y teniendo que $\mu^\circ(A) \geq \mu^\circ(\emptyset) = 0$, $\mu^\circ$ sobre $\mathcal{A}^\circ$ es una función aditiva y positiva. Para completar la prueba, veamos que $\mu^\circ$ es $\sigma$-aditiva. Sean $A_n \in \mathcal{A}^\circ$, $A = \displaystyle \sum A_n$. Tomamos los conjuntos $B_n = \displaystyle \sum_{k=1}^n A_k \in \mathcal{A}^\circ$. Tenemos que:
$$ \mu^\circ(D) \geq \mu^\circ (B_nD) + \mu^\circ (B_n^cD) \geq \sum_{k=1}^n\mu^\circ(A_kD) + \mu^\circ(A^cD) $$
Haciendo $n \to +\infty$
$$\mu^\circ(D) \geq \sum_{n=1}^{+\infty} \mu^\circ(A_nD) + \mu^\circ(A^cD) \geq \mu^\circ(AD) + \mu^\circ(A^cD)$$
Además queda demostrado que $A \in \mathcal{A}^\circ$, así que tenemos que $\mathcal{A}^\circ$ es un $\sigma$-campo.
Por último, tomando $D = A$, tenemos que $\mu^\circ(A) \geq \displaystyle \sum \mu^\circ(A_n) \Rightarrow \mu^\circ(A) = \sum \mu^\circ(A_n) \Rightarrow \mu^\circ$ es una medida.


\end{proof}

\begin{theorem}
de extensión de Carathéodory\\
Dada una medida $\mu$ sobre un álgebra $\mathcal{C}$, existe una extensión $\mu^\circ$ sobre la $\sigma$-álgebra minimal sobre $\mathcal{C}$ ($\mathcal{A}(\mathcal{C})$). Además, si $\mu$ es finita, la extensión es única.
\end{theorem}

\begin{proof}
  \

$\forall A \in \mathcal{C}, \forall D \in \mathcal{C}, \forall \varepsilon>0, \exists \{A_j\} \subset \mathcal{C}$ recubrimiento de $D$ tal que
$$\mu^\circ(D) + \varepsilon \geq \sum \mu(A_j) \geq \sum \mu (AA_j) + \sum \mu (A^cA_j) \geq \mu^\circ(AD) + \mu^\circ(A^cD) \Rightarrow$$
$$ \stackrel{\varepsilon \to 0}{\Rightarrow} \mu^\circ(D) \geq \mu^\circ(AD) + \mu^\circ(A^cD), A \in \mathcal{A}^\circ \Rightarrow \mathcal{C} \subseteq \mathcal{A}^\circ, \mathcal{A}(\mathcal{C}) \subseteq \mathcal{A}^\circ $$
Por los dos teoremas que hemos demostrado antes, la restricción $\mu^\circ$ sobre $\mathcal{A}(\mathcal{C})$ es una extensión de $\mu$ a una medida sobre $\mathcal{A}(\mathcal{C})$. Ahora, veamos que es única. Para ello, sean $\mu_1, \mu_2$ extensiones de $\mu$. Sea entonces el conjunto $\mathcal{M} = \{A: \mu_1(A) = \mu_2(A)\} \Rightarrow \mathcal{C} \subset \mathcal{M}$. Tenemos entonces que:
$$ \Omega \in \mathcal{C} \Rightarrow \Omega \in \mathcal{M} \Rightarrow \mu_1(\Omega) = \mu_2(\Omega) = \mu(\Omega) < \infty \Rightarrow \mu_1, \mu_2 \: finitos $$
$$ A \in \mathcal{M} \Rightarrow \mu_1(A^c) = \mu(\Omega) - \mu_1(A) = \mu(\Omega) - \mu_2(A) = \mu_2(A^c) \Rightarrow A^c \in \mathcal{M} $$
$$ A,B \in \mathcal{M} \Rightarrow \mu_1(A+B) = \mu_1(A) + \mu_1(B) = \mu_2(A) + \mu_2(B) = \mu_2(A+B) \Rightarrow A+B \in \mathcal{M}$$
Entonces, de momento $\mathcal{M}$ es un campo.\\

Sean ahora $\{A_n\} \subset \mathcal{M}$ monótona, entonces $\mu_1(\lim A_n) = \lim \mu_1(A_n) = \lim \mu_2(A_n) = \mu_2(\lim A_n) \Rightarrow \mathcal{M}$ cerrada frente al límite de sucesiones monótonas $\Rightarrow \mathcal{M} \sigma$-álgebra que contiene a $\mathcal{C}$, por tanto, contiene a la $\sigma$-álgebra minimal sobre $\mathcal{C}$
$$ \mathcal{A}(\mathcal{C}) \subseteq \mathcal{M} \Rightarrow \mu_1 = \mu_2 \: en \: \mathcal{A}(\mathcal{C})$$


\end{proof}

\chapter{Funciones medibles}

Antes de nada, vamos a puntualizar unas características básicas de la inversa de una función que podrían pasar desapercibidas:

\begin{lemma}
  Sea $X: \Omega \to \overline{\mathbb{R}}$ y $S',S, S_t \subset \mathbb{R} \forall t \in \mathbb{N}$
  \begin{itemize}
    \item $X^{-1}(S'-S)=X^{-1}(S')-X^{-1}(S))$  
    \item $X^{-1}(\cup_t S_t)=\cup_t X^{-1}(S_t)$
    \item $X^{-1}(\cap_t S_t)=\cap_t X^{-1}(S_t)$
  \end{itemize}
\end{lemma}

\begin{lemma}
  Sea $X: \Omega \to \overline{\mathbb{R}}$
  \begin{itemize}
    \item La imagen inversa de un $\sigma$-campo sobre $\mathbb{R}$ es un $\sigma$-campo
    \item La imagen inversa de un de $\mathcal{B}(\mathcal{C})$($\sigma$-campo de borel sobre $\mathcal{C}$) es el mínimo $\sigma$-campo sobre $X^{-1}(\mathcal{C})$.
    \item La clase de todos los conjuntos cuya antiimagen está en un $\sigma$-campo es un $\sigma$-campo.
  \end{itemize}
\end{lemma}

Una vez definidos los conceptos sobre espacios de medida y espacios probabilísticos, vamos a aproximarnos al concepto de función medible. Para ello, daremos dos definiciones de función medible, para demostrar más adelante que estas dos definiciones son equivalentes. Empecemos con la definición constructiva de función medible. Dado que nos interesa que el codominio sea $\mathbb{R}$, trabajaremos con funciones definidas entre un espacio medible ($\Omega$, $\mathcal{A}$) y ($\mathbb{R}$, $\mathcal{B}$), donde $\mathcal{B}$ representa la $\sigma$-álgebra de Borel.\\

Primero, se define la función puntual $X: \Omega \to \mathbb{R}$, que asigna a cada $\omega \in \Omega$ un número $x = X(\omega)$ único. Llamaremos a esta función variable aleatoria. Utilizaremos de aquí en adelante la siguiente notación:

\begin{itemize}
\item $[X]$ = dominio de la variable aleatoria $X$.
\item $[ X \leq Y] = \{\omega: X(\omega) \leq Y(\omega)\}$
\item $[ X = x] = \{\omega: X(\omega)= x\}$
\end{itemize}

Sea entonces la función $X = \displaystyle \sum_j x_jI_{A_j}$, donde $A_j$ son conjuntos medibles y $I_{A_j}$ denota la función indicadora de dicho conjunto. Estas funciones se llaman funciones elementales, y cuando el número de valores distintos que toma la función $X$ es finito, se conocen como funciones simples. Entonces, la definición constructiva de función medible es la que sigue:\\

\textbf{Definición constructiva de función medible:} Una función es medible si es límite de una sucesión de funciones simples $\{X_n\}$ convergentes.\\

Esta definición que hemos dado es constructiva, y por tanto, nos será muy útil para la definición constructiva de las integrales. No obstante, para enunciar y demostrar las propiedades de las funciones medibles, suele ser más útil la definición descriptiva siguiente:\\

\textbf{Definición descriptiva de función medible:} Sea una función $\varphi : \mathcal{A} \rightarrow \mathcal{B}$. Se dice que $\varphi$ es medible si $\forall B \in \mathcal{B} \Rightarrow \varphi^{-1}(B) \in \mathcal{A}$. Es decir, una función se dice medible si la imagen inversa de todo conjunto medible es medible.\\

No obstante, se puede dar, a raíz de esta, otra definición más económica:\\

Para la definición anterior, es suficiente con exigir la medibilidad de las imágenes inversas de los elementos de una subclase $\alpha$ para la cual la $\sigma$-álgebra minimal sobre $\alpha$ sea $\mathcal{B}$\\

Veamos ahora que las dos definiciones que hemos dado son equivalentes.

\begin{theorem} Teorema de medibilidad\\
Las definiciones constructiva y descriptiva de una función medible son equivalentes, y la clase de funciones medibles es cerrada bajo las operaciones usuales del análisis.
\end{theorem}

\begin{proof}
Sean $X_n$ funciones medibles en el sentido descriptivo. Entonces todos los conjuntos de la forma
$$[\inf X_n < x] = \cup [X_n < x], [-X_n < x] = [X_n > -x]$$
son medibles, y por tanto, las funciones
$$\sup X_n = - \inf (-X_n), \lim \inf (X_n) = sup (\inf_{k \geq n} X_k)$$
$$ \lim \sup X_n = - \lim \inf (-X_n)$$
son medibles en el sentido descriptivo. Por un lado, las funciones de este tipo son claramente cerradas bajo las operaciones de supremo, ínfimo, y límites. Además, toda función simple es medible en el sentido descriptivo, ya que todos los conjuntos $[X \leq x] = \displaystyle \sum_{x_j \leq x} A_j$ son medibles. Entonces, el límite de sucesiones convergentes de funciones simples es medible, y por tanto las funciones medibles en sentido constructivo lo son en sentido descriptivo.\\

Veamos la otra implicación, es decir, que las funciones medibles en sentido descriptivo lo son en sentido constructivo. Sea una función $X$ medible en sentido descriptivo. Entonces, las funciones
$$X_n = -nI_{[X < n]} + \sum_{k=-n2^n+1}^{n2^n}\frac{k-1}{2^n}I_{\big[\frac{k-1}{2^n} \leq X < \frac{k}{2^n} \big]} + nI_{X \geq n}, n \in \mathbb{N}$$
son simples. Entonces, dado que
$$\mid X_n(\omega) - X(\omega)\mid < \frac{1}{2^n} \quad para \quad \mid X(\omega)\mid < n$$
y
$$X_n(\omega) = \pm n \quad para \quad  X(\omega) = \pm \infty$$
se tiene entonces que $X_n \to X$ y esto, con lo anterior, prueba la equivalencia de las dos definiciones de función medible.


\end{proof}

%Este lemma era la definición descriptiva de función medible.

\begin{theorem}
  \ 
  
Sea $X:\Omega \longrightarrow \mathbb{R}$ medible y $g: \mathbb{R} \longrightarrow \mathbb{R}$ continua. Entonces, $ g(X)$ es medible
\end{theorem}

\begin{proof}
Si $X$ es elemental, tenemos que $\displaystyle g(X) = g\Big(\sum_jx_jI_{A_j}\Big)= \sum_j g(x_j)I_{A_j}$ función elemental, y por tanto, medible.\\

Si $X$ es el límite de funciones elementales:
$$g(X) = g(\lim X_n) = \lim g(X_n) = \lim \sum_j g(x_j)I_{A_j}$$


\end{proof}

Para este teorema introduceremos el concepto de función boreliana. Esta es una función $g:(\mathbb{R},\mathcal{B})\to(\mathbb{R},\mathcal{B})$ tal que $g^{-1}(B)\in \mathcal{B}, \forall B\in\mathcal{B}$

\begin{theorem}
  \
  
Sean $X$ función medible y $g$ función boreliana. Entonces, $g(X)$ es medible.
\end{theorem}

\begin{proof}
$$(g(X))^{-1}(B)=X^{-1}(g^{-1}(B))\in \mathcal{A}$$

\end{proof}

\section{Convergencia en probabilidad y convergencia casi segura}

Vamos a establecer ahora criterios que nos permitan comparar variables aleatorias. Dado un espacio de probabilidad ($\Omega$, $\mathcal{A}$, $\mathcal{P}$), se dice que dos variables aleatorias son equivalentes si:
$$ X\mathcal{R}Y \Leftrightarrow \mathcal{P}[X = Y] = 1$$
Análogamente:
$$X(\omega) = Y(\omega) \hspace{1cm} \forall \omega \in \Omega \setminus \Lambda, \  \mathcal{P}(\Lambda) = 0$$

\subsection{Definición de convergencia en probabilidad}

Si $P[|X_n- X| \geq \varepsilon]\rightarrow 0$, entonces se dice que $X_n\stackrel{\mathbb{P}}{\longrightarrow} X$, es decir, $X_n$ converge en probabilidad a $X$\\

\begin{lemma}
$X_n\stackrel{\mathcal{P}}{\longrightarrow} X \wedge X_n\stackrel{\mathcal{P}}{\longrightarrow} Y \Longrightarrow X\mathcal{R}Y$\\
\end{lemma}

\begin{proof}
$$|X-Y |=|X - X_n - Y + X_n | \leq |X_n - X|+| X_n-Y |$$

$$P[|X - Y| \geq \varepsilon] \leq P[|X_n - X | \geq \frac{\varepsilon}{2} ] + P[| X_n - Y | \geq \frac{\varepsilon}{2}] \rightarrow 0$$
Y por tanto, $X \mathcal{R} Y$


\end{proof}

\subsection{Definición de convergencia uniforme en u}

Se dice que una variable aleatoria $X$ converge uniformemente en $u$ si se cumple que $\mathcal{P}[\mid X_{n+u} - X_n \mid \geq \varepsilon] \to 0$

\begin{lemma}
\

La convergencia en probabilidad implica la convergencia uniforme
\end{lemma}

\begin{proof}
$$ \mid X_{n+u} - X_n \mid \ \leq \ \mid X_{n+u} - X \mid + \mid X - X_n \mid \Rightarrow$$
$$ \Rightarrow \mathcal{P}[\mid X_{n+u} - X_n \mid \geq \varepsilon] \ \leq \ \mathcal{P}\Big[\mid X_{n+u} - X \mid \geq \frac{\varepsilon}{2}\Big] + \mathcal{P}\Big[\mid X_{n} - X \mid \geq \frac{\varepsilon}{2}\Big] \to 0$$


\end{proof}

\subsection{Propiedades}

\begin{enumerate}
\item $ X_n \stackrel{\mathcal{P}}{\longrightarrow} X \wedge Y_n \stackrel{\mathcal{P}}{\longrightarrow} Y \Longrightarrow X_n + Y_n \stackrel{\mathcal{P}}{\longrightarrow} X + Y$
\item $X_n \stackrel{\mathcal{P}}{\longrightarrow} X \Longrightarrow KX_n \stackrel{\mathcal{P}}{\longrightarrow} KX$
\item $X_n \stackrel{\mathcal{P}}{\longrightarrow} K$ (es decir, que degenera), entonces ${X_n}^2 \stackrel{\mathcal{P}}{\longrightarrow} K^2$.\\
Para demostrarlo basta notar que: $ {X_n}^2 - K^2 = (X_n + K) (X_n -K)$
\item $X_n \stackrel{\mathcal{P}}{\longrightarrow} a \wedge Y _n \stackrel{\mathcal{P}}{\longrightarrow} b \Longrightarrow X_nY_n \stackrel{\mathcal{P}}{\longrightarrow} a \cdot b$\\

Para demostrarlo tenemos que notar que:
$$X_n Y_n = \dfrac{(X_n + Y_n)^2-(X_n - Y_n)^2}{4}\stackrel{\mathcal{P}}{\longrightarrow} \dfrac{(a+b)^2-(a-b)^2}{4}=ab$$

\item $\forall n\in\mathbb{N}$ con $X_n\not = 0, \ X_n \stackrel{\mathcal{P}}{\longrightarrow} 1 \Longrightarrow \dfrac{1}{X_n} \stackrel{\mathcal{P}}{\longrightarrow} 1$
\item $\forall n\in\mathbb{N}$ con $Y_n\not = 0, \ X_n \stackrel{\mathcal{P}}{\longrightarrow} a \wedge Y_n \stackrel{\mathcal{P}}{\longrightarrow} b\not =0 \Longrightarrow X_nY_n^{-1}\stackrel{\mathcal{P}}{\longrightarrow} ab^{-1}$
\item $X_n \stackrel{\mathcal{P}}{\longrightarrow} X \wedge Y_n\stackrel{\mathcal{P}}{\longrightarrow} Y \Longrightarrow X_nY_n\stackrel{\mathcal{P}}{\longrightarrow} XY$\\

Para demostrarlo basta notar que: $(X_n - X)(Y_n - Y)\stackrel{\mathcal{P}}{\longrightarrow} 0$ y luego $X_nY_n-XY_n-X_nY+XY\stackrel{\mathcal{P}}{\longrightarrow} 0$ (Por la propiedad siguiente)
\item $X_n\stackrel{\mathcal{P}}{\longrightarrow} X$, entonces $YX_n\stackrel{\mathcal{P}}{\longrightarrow} YX$
\end{enumerate}

\begin{theorem}
Si $X_n\stackrel{\mathcal{P}}{\longrightarrow} X$ y $g(\cdot)$ es continua, entonces se cumple que:
$$g(X_n)\stackrel{\mathcal{P}}{\longrightarrow} g(X)$$
\end{theorem}

\begin{proof}
$\forall \varepsilon > 0, \exists k>0 : \mathcal{P}[|X| \geq k] < \frac{\varepsilon}{2}$. Consideramos el compacto $[-k,k], \ A = [|X| < k], \ B_n = [|X_n - X| < \delta] \ y \ C = [|g(X_n) - g(X)| < \varepsilon]$. Tenemos que ver si $\mathcal{P}(B_n)\to 1 \Rightarrow \mathcal{P}(C)\to 1 $:
$$ AB \subset C \Rightarrow C^c \subset A^c \cup B^c \Rightarrow \mathcal{P}(C^c) \leq \mathcal{P}(A^c) + \mathcal{P}(B^c) \leq $$
$$ \leq \mathcal{P}[|X| \geq k] + \mathcal{P}[|X_n - X| \geq \delta] \leq \frac{\varepsilon}{2} + \frac{\varepsilon}{2} = \varepsilon$$
\end{proof}

\subsection{Definición de convergencia casi segura}

Una sucesión de variables aleatorias, ${ X_n }$, converge con probabilidad 1, o de forma casi segura, a una variable aleatoria $X$ (que puede degenerar en una constante K) cuando se cumple que:
$$P(\lim_{n\rightarrow\infty}X_n=X)=1$$
De esta forma interpretamos que $X_n\stackrel{c.s}{\longrightarrow}X$ cuando la probabilidad de que en el límite la sucesión de variables  aleatorias y aquella a la que converge sean iguales es $1$.

\begin{theorem}
$$X_n\stackrel{c.s}{\longrightarrow}X\Longrightarrow X_n\stackrel{\mathcal{P}}{\longrightarrow}X$$
$$X_n \text{ finitas c.s.,  }
X_n\stackrel{\mathcal{P}}{\longrightarrow}X
\Longrightarrow 
\exists X_{n_k}\,:\, X_{n_k}
\stackrel{c.s}{\longrightarrow}_{k\rightarrow\infty}X$$
\end{theorem}
\begin{proof}
$$0=\displaystyle\lim_{n\rightarrow\infty}P\bigcup_{m\geq n}[|X_m - X|\geq \epsilon ]\geq\displaystyle\lim_{n\rightarrow\infty}P[|X_n - X|\geq \epsilon ]$$


\end{proof}

\begin{lemma}
Sea $\mathbb{X}_n$ finita. Entonces se tiene que:
\\
$$\mathbb{X}_n \stackrel{c.s.}{\longrightarrow}\mathbb{X}\Longleftrightarrow \mathbb{X}_n \text{ es de Cauchy}$$
\end{lemma}
\begin{proof}
$$\Longrightarrow$$
$\exists n_0\, : \, \forall n\geq n_0$ se tiene que:
$$|\mathbb{X}_n(\omega)-\mathbb{X}(\omega)|<\frac{\epsilon}{2}
\hspace{0.5cm}  \omega \in A,\ P(A)=1$$
$$|\mathbb{X}_n(\omega)-\mathbb{X}_m(\omega)|\leq|\mathbb{X}_n(\omega)-\mathbb{X}(\omega)|+|\mathbb{X}_m(\omega)-\mathbb{X}(\omega)|<\epsilon$$
$$\Longleftarrow$$
\begin{center}
${\mathbb{X}_n(\omega)}$ de Cauchy
\end{center}
$$\lim\mathbb{X}_n(\omega)=\mathbb{X}(\omega)\hspace{0.5cm} \omega \in A$$

\end{proof}

\chapter{Esperanza matemática}

\section{Integral de una función de conjunto}

Para el cálculo de probabilidades, nos será muy útil el concepto de integral sobre funciones de conjunto. Este concepto nos servirá para calcular lo que se conoce como la esperanza matemática de una variable aleatoria $X$. En este apartado trabajaremos sobre el espacio de probabilidad $(\Omega, \mathcal{A}, \mathcal{P})$. Comenzaremos definiendo la integral para las funciones simples, para dar luego una definición de integral para funciones no negativas y por último para funciones cualesquiera.\\

Sea entonces $\{A_k\} \in \mathcal{A}$, tal que $\displaystyle \sum_k A_k = \Omega$, partición medible del espacio. Sea entonces la función simple $X = \displaystyle \sum_{j=1}^m x_jI_{A_j},\ x_j \geq 0$. La integral de la función $X$ se define como:
$$\int_{\Omega} X d\mathcal{P} = \sum_{j=1}^m x_j\mathcal{P}_{A_j}$$
Ahora, para cualquier función no negativa $X$, se define la integral de la función como:
$$\int_{\Omega} X d\mathcal{P} = \lim \int_{\Omega}X_n d \mathcal{P}$$
Donde $\{X_{n}\} \to X$. Finalmente, la integral en $\Omega$ de una función medible $X$ se define como:
$$\int_{\Omega} X d\mathcal{P} = \int_{\Omega}X^{+} d\mathcal{P} - \int_{\Omega} X^{-} d\mathcal{P}$$
donde $X^{+} = XI_{[X\geq0]}$ y $X^{-} = -XI_{[X<0]}$. Si $\displaystyle \int_{\Omega}Xd\mathcal{P}$ es finita, es decir, si los dos términos de la diferencia anterior son finitos, entonces se dice que $X$ es integrable en $\Omega$. Ahora, una vez definida la integral, vamos a ver algunas de sus propiedades. Tenemos primero una serie de propiedades relacionadas con la aditividad de la integral. Sean $X,Y$ dos funciones medibles, entonces (no se demostrarán las propiedades triviales):\\

$\displaystyle \int (X+Y)d\mathcal{P} = \int Xd\mathcal{P} + \int Yd\mathcal{P}$\\

\begin{proof}

Sean $X = \displaystyle \sum x_jI_{A_j}$ e $Y = \displaystyle \sum y_kI_{B_k}$. Entonces $X+Y = \displaystyle \sum_k \sum_j x_j I_{A_jB_k} + \sum_k \sum_j y_k I_{A_jB_k} = \sum_k \sum_j (x_k+y_j) I_{A_jB_k}$\\

Si calculamos ahora las integrales:
$$ \int (X+Y)d\mathcal{P} = \sum_k + \sum_j (x_j + y_k) \mathcal{P}(A_jB_k) =$$
$$ = \sum_j x_j \mathcal{P}(A_j) + \sum_k y_k \mathcal{P}(B_k) = \int Xd\mathcal{P} + \int Yd\mathcal{P}\\$$

\end{proof}

$\displaystyle \int_{A+B} X d\mathcal{P} = \int_A X d\mathcal{P} + \int_B X d \mathcal{P}$\\

$\displaystyle \int cXd\mathcal{P} = c\int Xd\mathcal{P}$\\

Veamos ahora algunas propiedades relacionadas con el orden:\\

$X \geq 0 \rightarrow \displaystyle \int X d\mathcal{P} \geq \int 0 = 0$\\

$X \geq Y \rightarrow \displaystyle \int X d\mathcal{P} \geq \int Y d\mathcal{P}$\\

$\displaystyle X \stackrel{c.s}{=} Y \rightarrow \int X d\mathcal{P} = \int Y d\mathcal{P}$\\

\section{Esperanza matemática}

\subsection{Definición general}

Dado un espacio probabilístico $(\Omega, \mathcal{A}, P)$ y una v.a. $X$, se define la esperanza matemática de $X$ como:
$$E(X)=\int_{\Omega}XdP=\int_{\omega \in \Omega}{X(\omega)dP(\omega)}$$

La esperanza matemática se define de forma distinta si estamos trabajando con una v.a. discreta o continua. En el caso de una v.a. discreta, sea $P(x_i)$ su función de probabilidad. Definimos la esperanza matemática como:
$$E(X)=x_1P[X=x_1]+\ldots + x_n P[X=x_n]=\sum_{i=1}^{n}{x_i P(x_i)}$$
En el caso de una v.a. continua, sea $f(x)$ la función de densidad. Definimos la esperanza matemática como:
$$E(X)=\int_{-\infty}^{+\infty} {xf(x)dx}$$

\subsection{Teorema de la convergencia monótona}

Una vez vista la definición de integral y esperanza matemática y algunas de sus propiedades, vamos a enunciar y demostrar un teorema de convergencia que nos será de mucha utilidad para el estudio de variables aleatorias. Veamos su enunciado y demostración:
\begin{theorem}
  \
  
Teorema de la convergencia monótona para funciones medibles no negativas\\

Sea $\{X_n\} \geq 0$ tal que $X_n \uparrow X$. Entonces, se tiene que $\displaystyle \int X_n \uparrow \int X$
\end{theorem}

\begin{proof}
\

Tomamos las sucesiones $X_{kn} \uparrow X_k$. La sucesión $Y_n = \displaystyle \max_{k \leq n} X_{kn}$ es una sucesión de funciones simples no negativas y no decreciente, y además
$$X_{kn} \leq Y_n \leq X_n 
\hspace{1cm}
\int X_{kn} \leq \int Y_n \leq \int X_n$$
Ahora, cuando hacemos $\lim_{n \rightarrow \infty}$, tenemos que
$$X_k \leq lim Y_n \leq X 
\hspace{1cm}
\int X_k \leq \int lim Y_n \leq \lim \int X_n$$
Por último, cuando $k \rightarrow \infty$, obtenemos
$$X \leq lim Y_n \leq X 
\hspace{1cm}
lim \int X_n \leq \int lim Y_n \leq lim \int X_n$$
De donde extraemos que $\lim Y_n = X$ y que $\displaystyle \int X = lim \int X_n$


\end{proof}

\subsection{Teorema de Fatou-Lebesgue}
Veamos ahora un resultado que nos da información sobre el límite superior e inferior de una sucesión de variables aleatorias, en caso de que estas estén acotadas por variables aleatorias integrables

\begin{theorem}
de Fatou-Lebesgue\\

Sean $Y$, $Z$ dos funciones integrables (pueden no mantener su signo) y $X_n$ una sucesión de variables aleatorias, entonces:
\begin{enumerate}
\item $\ \forall n\in\mathbb{N} \ Y \leq X_n \Longrightarrow$ $E(\lim \inf X_n) \leq \lim \inf E(X_n)$
\item $\ \forall n\in\mathbb{N}\ X_n \leq Z \Longrightarrow E(\lim \sup X_n)\geq \lim \sup E(X_n)$
\end{enumerate}
\end{theorem}

\begin{proof}
Si $X_n \geq 0$, tenemos que:
$$ X_n \geq Y_n = \inf_{k \geq n} X_n \uparrow \lim \inf X_n \stackrel{TCM}{\implies} \lim \inf E(X_n) \geq \lim Y_n = E(\lim \inf X_n)$$
En otro caso:
\begin{enumerate}
\item $X_n - Y \geq 0 \Rightarrow \lim \inf E(X_n) - \lim \inf E(Y) \geq E(\lim \inf X_n) - E(\lim \inf Y)$
\item $Z - X_n \geq 0 \Rightarrow - \lim \inf E(-X_n) \geq - E(\lim \inf - X_n) \Rightarrow \lim \sup E(X_n) \leq E(\lim \sup X_n)$
\end{enumerate}

\end{proof}

\subsection{Teorema de la convergencia dominada}
\begin{theorem}
Sea $X_n$ una sucesión de variables aleatorias t.q.  $X_n \stackrel{\mathcal{P}}{\rightarrow} X$ o $X_n \stackrel{c.s.}{\rightarrow} X$  . Si $\exists Y \, t.q. \, |X_n |\leq Y$. Entonces se tiene que $X$ es integrable y:
$$\lim\int X_n\, dP=\int X\, dP$$
\end{theorem}

\begin{proof}
Por un lado, si tenemos la convergencia casi segura, $|X_n| \leq Y \Rightarrow -Y \leq X_n \leq Y \stackrel{F-L}{\rightarrow} \displaystyle \int X_n \to \int X$.
Por otro, si la convergencia es en probabilidad, nos bastaría con probar que si $Y_n \stackrel{\mathcal{P}}{\rightarrow} 0 \Rightarrow \displaystyle \int Y_n \stackrel{\mathcal{P}}{\rightarrow} 0$. Tomamos entonces $Y_n = |X_n - X|$
$$ Y_n \stackrel{\mathcal{P}}{\rightarrow} 0 \Rightarrow \exists \{Y_n'\} \subset \{Y_n\} \, t.q. \, \int Y_n' \rightarrow \limsup \int Y_n$$
De nuevo, tenemos que $Y_n' \stackrel{\mathcal{P}}{\rightarrow} 0$. Tomando ahora $\{Y_n''\} \subset \{Y_n'\} \, t.q. \, Y_n'' \stackrel{c.s.}{\rightarrow} 0 
\implies 
\displaystyle \int Y_n'' \rightarrow 0 \stackrel{unic.lim.}{\implies} \lim \sup \int Y_n = 0 \stackrel{Y_n \geq 0}{\implies} \int Y_n \rightarrow 0$


\end{proof}

\subsection{Desigualdades}

Vamos a ver algunas desigualdades que nos permitirán a posteriori simplificar el cálculo de carácterísticas asociadas a variables aleatorias.\\

\textbf{Definición: Momento r-ésimo}\\

Dado un espacio probabilístico ($\Omega$, $\mathcal{A}$, $\mathcal{P}$) y una variable aleatoria $X$, se define el momento r-ésimo o de orden r de la variable $X$ a:
$$ EX^r = \int X^r \, d \mathcal{P} $$
Dicho momento existe si $E|X|^r < \infty$

\subsubsection{Desigualdad $C_r$}
\begin{lemma}

\
  
$|a+b|^r\leq C_r |a|^r + C_r |b|^r$ con $C_r=1$ si $r\leq1$ y $C_r=2^{r-1}$ si $r>1$\\

Entonces se tiene $E|X+Y|^r \leq C_r E|X|^r + C_r E|Y|^r$
\end{lemma}

\subsubsection{Desigualdad de Hölder}
$E|XY| \leq \left(E|X|^r \right)^{\frac{1}{r}}  \left(E|Y|^s \right)^{\frac{1}{s}}$ con $r>1$ ; $ \dfrac{1}{r}+\dfrac{1}{s}=1$

\subsubsection{Caso particular: Desigualdad de Schwarz}
$E|XY|\leq \left(E|X|^2 \right)^{\frac{1}{2}} \left(E|Y|^2 \right)^{\frac{1}{2}}$

\subsubsection{Desigualdad de Minkowsky}
Si $r \geq 1 \text{, entonces } (E|X+X'|^r)^{\frac{1}{r}} \leq \left(E|X|^r \right)^{\frac{1}{r}} + \left(E|X'|^r \right)^{\frac{1}{r}}$

\subsubsection{Desigualdad básica}
Sea $X$ una v.a. y $g$ una función boreliana no negativa en $\mathbb{R}$, se tiene:

\begin{description}
\item[i)] Si $g$ es par y no decreciente en $[0,\infty)$, entonces $\forall a>0$:
$$\dfrac{Eg(X)-g(a)}{c.s. \sup\, g(X)}\leq P[|X|\geq a]\leq \dfrac{E g(X)}{g(a)}$$
\item[ii)] (Desigualdad de Tchevychev) 

Si $g(x)=x^2$ y cambiamos la v.a. por $X-EX$, por ser $g$ no decreciente en $\mathbb{R}$, podemos prescindir del valor absoluto:
$$ P[X-EX \geq a] \leq \frac{Var X}{a^2}$$
\item[iii)] (Desigualdad de Markov) Si $g(X) = |X|^r$
$$ \frac{EX^r - a^r}{c.s.\sup X^r} \leq P[|X| \geq a] \leq \frac{EX^r}{a^r} $$
\end{description}

Vamos a demostrar que en efecto se cumple la desigualdad básica:
\begin{proof}
Definimos $A = [|X| \geq a]$
$$ Eg(A) = \int_A g(X) d\mathcal{P} + \int_{A^c} g(X) d\mathcal{P} $$
$$ g(a)\mathcal{P}(A)\ \leq\ \int_A g(X)d\mathcal{P} \ \leq \ c.s.\sup g(X)\mathcal{P}(A) $$
$$ 0 \ \leq \ \int_{A^c} g(X) d\mathcal{P} \ \leq \ g(a) $$
$$ g(a)\mathcal{P}(A) \ \leq\ Eg(X) \ \leq \ c.s.\sup g(X)\mathcal{P}(A) + g(a) $$

\end{proof}

\subsubsection{Aplicaciones}
\begin{enumerate}
\item Usando la desigualdad básica con $g(X) = X^r$ y $E|X_n - X|^r \to 0 \Rightarrow X_n \stackrel{\mathcal{P}}{\rightarrow} X$
\item ($\Omega$, $\mathcal{A}$, $\mathcal{P}$) $\stackrel{\mathcal{P}}{\rightarrow}$ ($\mathbb{R}$, $\mathcal{B}$), definimos ahora $\mathcal{P}^X(\mathcal{B}) = \mathcal{P}(X^{-1}(\mathcal{B}))$. $\mathcal{P}^X$ es una probabilidad, ya que $\mathcal{P}^X(\mathbb{R}) = 1, \mathcal{P}^X(\mathcal{B}) \geq 0, \mathcal{P}^X(\sum B_j) = \sum \mathcal{P}^X(B_j)$, y $EX = \displaystyle \int_{x \in \mathbb{R}} x d \mathcal{P}^X(x)$; $Eg(X) = \displaystyle \int_{y \in \mathbb{R}} g d \mathcal{P}^Y(y) = \displaystyle \int_{x \in \mathbb{R}} g(x) d \mathcal{P}^X(x)$
\end{enumerate}

\section{Convergencia en r-medida}
Una sucesión de variables aleatorias $X_n$ converge en r-medida a $X$ si:
$$ E|X_n - X|^r \to 0$$
Lo notamos como $X_n \stackrel{r}{\rightarrow} X$

\begin{lemma}
$X_n \stackrel{r}{\rightarrow} X \Rightarrow E|X_n|^r \rightarrow E|X|^r$
\end{lemma}

\begin{proof}
Tenemos dos casos, para $r \leq 1$ y $r \geq 1$\\

Para $r \leq 1$: $|E|X_n|^r - E|X|^r| \leq E|X_n - X|^r \rightarrow 0$\\

Para $r \geq 1$: $|(E|X_n|^r)^{\frac{1}{r}} - (E|X|^r)^{\frac{1}{r}}| \leq (E|X_n - X|^r)^{\frac{1}{r}} \rightarrow 0$
\end{proof}

\section{Función de distribución, distribución inducida y función generatriz de momentos}
\subsection{Función de distribución}
Función puntual que se define sobre la recta real, no negativa, no decreciente, continua por la izquierda y que cumple que $F(-\infty) = 0$, $F(+\infty) = 1$\\

\subsection{Distribución inducida}
Se conoce como distribución inducida por la variable aleatoria $X$ a ($\mathbb{R}$, $\mathcal{B}$, $\mathcal{P}^X$), donde la medida de probabilidad $\mathcal{P}^X$ se define como:
$$ \mathcal{P}^X(B) = \mathcal{P}(X^{-1}(B)) = \mathcal{P}[X\in B], \, \forall B \in \mathcal{B} $$
Tenemos entonces que:
\begin{itemize}
\item $\mathcal{P}^X(B) \geq 0,\  \forall B \in \mathcal{B}$
\item $\displaystyle \mathcal{P}^X(\sum_j B_j) = \mathcal{P}(X^{-1}(\sum_jB_j)) = \mathcal{P}(\sum_j X^{-1}(B_j)) = \sum_j \mathcal{P}(X^{-1}(B_j)) = \sum_j \mathcal{P}^X(B_j)$
\item $\displaystyle EX = \int_{\omega \in \Omega} X(\omega) d \mathcal{P(\omega)} = \int_{x \in \mathbb{R}} x d\mathcal{P}^X(x)$
\end{itemize}

\subsection{Función generatriz de momentos}
Llamamos función generatriz de momentos de la variable aleatoria $X$ a la función:
$$ M_X(t) = E(e^{tX}) \hspace{1cm} t \in \mathbb{R} $$
\subsection{Función característica}
Dada una variable aleatoria $X$ continua, definimos su función característica:
$$\varphi_X(t)=E(e^{itX})=\int_{-\infty}^{\infty}e^{itx}f_{X}(x)=E(cos(tX))+iE(sen(tX))$$

\chapter{Independencia}
\section{Concepto de independencia}
\textbf{Definición: }Sea un espacio probabilístico $(\Omega,\mathcal{A},\mathcal{P})$. Se dice que los eventos $\{A_t : t\in T \}\in \mathcal{A}$ son independientes si para todo subconjunto finito $(t_1,\ldots ,t_n)\in T$ se tiene que:
$$P\bigcap_{k=1}^n A_{t_k}=\displaystyle\prod_{k=1}^{n}P\, A_{t_k}$$
De hecho, el concepto de independencia es relativo a la familia de clases. Una colección de clases de eventos se dice independiente si seleccionando un evento arbitrario de cada clase dichos eventos son independientes. 

\begin{theorem}
  \
  
Se tiene que las subclases de clases independientes son independientes.
\end{theorem}

\begin{proof}
	La demostración es trivial puesto que si todos los conjuntos de una clase cumplen la propiedad, también la han de cumplir un subconjunto de esta.
\end{proof}

\section{Suma de variables aleatorias independientes}
\subsubsection{Teorema de las funciones Borel}
\begin{theorem}
Las funciones de Borel $g(X_i)$ de funciones aleatorias independientes son independientes.
\end{theorem}


\begin{proof}
Las clases $\mathcal{B}_t\subset \mathcal{B}(X_t) = X_t'(\mathcal{B})$ son indep. $\implies$ g induce un sub $\sigma$-campo $\mathcal{B}_t$ de eventos contenidos en $\mathcal{B}(X_t) \ (\mathcal{B}(g(X_t)) \subset \mathcal{B}(X_t))$
\\


($X_t$ indep.)

Sea $g$ función de Borel $\implies$ $(g \circ X_t)$ son independientes.

\end{proof}

\subsubsection{Teorema de extensión}
\begin{theorem}
Los $\sigma$-campos minimales sobre clases independientes cerradas para la intersección finita son independientes.
\end{theorem}
Esto tiene tres aplicaciones fundamentales:
\begin{enumerate}
\item Si los eventos $A_t$ son independientes, también lo son los $\sigma$ campos $(A_t , A_t^c , \emptyset , \Omega)$.
\item Si la imagen inversa $\mathcal{C}_t$ de las clases de todos los intervalos $(-\infty, x_t)$ en los espacios de Borel $R_t$ son independientes, también lo son las imágenes inversas $\mathcal{B}_t$ de los campos de Borel en $R_t$. Para todo $\mathcal{C}_t$, éste es cerrado para la intersección finita y $\mathcal{B}_t$ es el $\sigma$ campo minimal sobre $\mathcal{C}_t$
\item Sean $\mathcal{B}_t$ campos de eventos y sea $T_s$ un subconjunto de $T$. El $\sigma$ campo $\mathcal{B}_T$ compuesto por componentes $\mathcal{B}_t$, con $t\in T_s$, es el $\sigma$ campo minimal sobre la clase $\mathcal{C}$ de todas las intersecciones finitas de eventos $A_t$ y contiene a todos sus componentes
\end{enumerate}
\begin{proof}

Sean $\mathcal{C}_1 \dots \mathcal{C}_n \ con \ \mathcal{C}_i^* \subset \mathcal{C}_i \ i=1\dots n \implies \mathcal{C}_i^* \ son \ independientes.$
\\
\begin{itemize}
\item $\o$ y los eventos c.s.

$\{\mathcal{C}_1 \cup \o \cup \Omega \cup \dots \}, \dots , \{\mathcal{C}_n \cup \o \cup \Omega \cup \dots \}$ $A_i' \subset A_i \in \mathcal{C}_i \ i=1\dots n.$
\\
\item Las diferencias apropiadas a sus elementos, en concreto los complementarios.

$P(A_1 - A_1') A_2\dots A_n = PA_1\dots A_n - PA_1'\dots A_n = \prod_{i=1}^{n}PA_{i} - PA_1'\prod_{i=1}^{n}PA_i$
\\
\item La suma contable de sus elementos.

$P(\sum_{j=1}^{n}A_1^j)A_2\dots A_n = \sum_{j=1}^{n}PA_1^jA_2\dots A_n = \sum_{j=1}^{n}PA_1^jPA_2\dots PA_n = P(\sum_{j=1}^{n}A_1^j)PA_2\dots PA_n$
\\
\item Los limites de sequencias de sus elementos.

$\lim A_{1}^m \rightarrow A_{1} \ si \ m \rightarrow \infty$

$PA_1A_2\dots A_n \leftarrow PA_1^mA_2\dots A_n = PA_1^mPA_2\dots PA_n \rightarrow PA_1PA_2\dots PA_n$
\end{itemize}

\end{proof}
\subsubsection{Teorema de composición}
\begin{theorem}
Los $\sigma$-campos compuestos son independientes si y solo si sus $\sigma$-subcampos finitamente compuestos son independientes.
\end{theorem}
\subsection{Familia de variables aleatorias}
Una familia de variables aleatorias $X_{T_s}=\{ X_t : t\in T_s \}$ induce $\sigma$-campos $\mathcal{B}(X_{T_s})$ de eventos.
\subsubsection{Teorema de las familias}
\begin{theorem}
Las familias de variables aleatorias son independientes si y solo si sus subfamilias finitas son mutuamente independientes.
\end{theorem}
\section{Multiplicación de variables aleatorias independientes}
\textbf{Definición: } Las variables aleatorias $X_t\, , \, t\in T$ son independientes si para cada clase finita $(S_{t_1},\ldots , S_{t_n})$ de conjuntos de Borel en $\mathbb{R}$,
$$P\bigcap_{k=1}^{n}[X_{t_k}\in S{t_k}]=\displaystyle\prod_{k=1}^n P[X_{t_k}\in S_{t_k}]$$

\subsubsection{Lema de la multiplicación}
\begin{lemma}
Si $X_1, \ldots , X_n$ son variables aleatorias independientes no negativas, entonces $E\displaystyle\prod_{k=1}^{n}X_k=\displaystyle\prod_{k=1}^{n} E X_k$
\end{lemma}
\begin{proof}
Basta demostrar el resultado para dos variables, dado que el resto se demuestra por inducción. En primer lugar, veamos que se verifica la condición para funciones simples.\\

Sean $X = \sum_j x_j I_{A_j}, Y = \sum_k y_k I_{B_k}$ variables simples no negativas, donde $A_j = [X = x_j], B_k = [Y = y_k]$.\\

Por ser $X,Y$ independientes, $P(A_jB_k) = P(A_j)P(B_k) \Rightarrow E(XY) = \sum_{k,j} P(A_jB_k) = \sum_{j,k} P(A_j)P(B_k) = E(X)E(Y)$\\

Sean ahora $X,Y$ variables aleatorias no negativas independientes. 
Su independencia implica la de
$$ A_{nj} = \left[\frac{j-1}{2^n} \leq X < \frac{j}{2^n}\right],\  B_{nk} = \left[\frac{k-1}{2^n} \leq Y < \frac{k}{2^n}\right]$$
Finalmente, tendremos la independencia de las variables aleatorias simples:
$$ X_n = \sum_{j=1}^{n2^n} \frac {j-1}{2^n}I_{A_{nj}}\hspace{0.5cm} Y_n = \sum_{k=1}^{n2^n} \frac {k-1}{2^n}I_{B_{nk}}$$
Tenemos entonces que:
$$ 0 \leq X_n \uparrow X, \ 0 \leq Y_n \uparrow Y \implies X_nY_n \uparrow XY$$
$$ E(X_nY_n) = E(X_n)E(Y_n)$$
De donde se extrae por el teorema de la convergencia monótona que $E(XY) = E(X)E(Y)$


\end{proof}

\subsubsection{Teorema de la multiplicación}
\begin{theorem}
Sean $X_1,\ldots, X_n$ variables aleatorias independientes. Si son integrables, entonces también lo es su producto con: $$E\displaystyle\prod_{k=1}^n X_k=\displaystyle\prod_{k=1}^n E X_k$$
\\\\
A la inversa, si su producto es integrable y ninguna de ellas degenera a $0$, entonces todas ellas son integrables.
\end{theorem}

\begin{proof}
$\newline$
$\Rightarrow$

La independencia de $X,Y$ nos da la independencia de $X', Y'$, donde $X' \in \{X^+, X^-, |X|\}$, $Y' \in \{Y^+, Y^-, |Y|\}$, por lo que $E(X'Y') = E(X')E(Y')$.\\

Si $X,Y$ son integrables, también lo son $X',Y'$ y por la anterior igualdad también lo es $X'Y'$, lo cual nos da también $|XY|$, y por tanto, $XY$ integrables con:
$$ E(XY) = E[(X^+-X^-)(Y^+-Y^-)] = E(X^+)E(Y^+) - E(X^+)E(Y^-) $$
$$ - E(X^-)E(Y^+) + E(X^-)E(Y^-) = E(X)E(Y) < \infty$$

$\Leftarrow$

Si $XY$ son integrables, entonces $E(|X|)E(|Y|) = E (|XY|)$

Como ni $X$ ni $Y$ degeneran a $0$, 
% por tanto $E(|X|)$ y $E(|Y|)$ no desvanecen % No sé lo que es desvanecer, Loeve facts
$E(|X|)$ y $E(|Y|)$ son finitas y, por tanto, $X$ e $Y$ son integrables.


\end{proof}

\subsubsection{Teorema de equivalencia}
\begin{theorem}
Para cada clase finita de conjuntos de Borel $S_t$ y de puntos $x_t, u_y \in \mathbb{R}$, son equivalentes las tres definiciones siguientes de independencia de variables aleatorias $X_t$:
\begin{enumerate}
\item $P\bigcap_{k=1}^n [X_{tk}\in S_{tk}]=\displaystyle\prod_{k=1}^n P[X_{tk}\in S_{tk}]$
\item $F_{t_1\ldots t_n}(x_{t_1},\ldots ,x_{t_n})=F_{t_1}(x_{t_1})\ldots F_{t_n}(x_{t_n})$
\item $f_{t_1\ldots t_n}(u_{t_1},\ldots ,u_{t_n})=f_{t_1}(u_{t_1})\ldots f_{t_n}(u_{t_n})$
\end{enumerate}
\end{theorem}
\begin{proof}$\,$\\
$i) \Rightarrow ii)$\\
Tomamos $S_t = ]-\infty, x_t[$ y aplicamos la definición.\\
$ii) \Rightarrow i)$\\
$S_t = ]-\infty, x_t[$,  $\mathcal{C}_t = {]-\infty, x_y[, t \in T}$ 
son independientes. 
Por el teorema de extensión, los $\sigma$-campos minimales sobre ellos son independientes.\\
$i) \Rightarrow iii)$\\
Aplicación directa del teorema de la multiplicación.\\
$iii) \Rightarrow ii)$\\
Tenemos que $F_{t_1,...,t_n}([a_{t_1},...,a_{t_n};b_{t_1},...,b_{t_n}]) = F_{t_1}([a_{t_1},...,a_{t_n};b_{t_1},...,b_{t_n}])$\\$F_{t_2}([a_{t_1},...,a_{t_n};b_{t_1},...,b_{t_n}])...F_{t_n}([a_{t_1},...,a_{t_n};b_{t_1},...,b_{t_n}])$. Haciendo $a_t \rightarrow - \infty$, $b \uparrow x_t$ tenemos $ii)$


\end{proof}

\subsubsection{Lema de Borel-Cantelli}
\begin{lemma}
Lema de Borel-Cantelli
\\\\
Sea $(\Omega,\mathcal{A},P) \, , \, A_n \in \mathcal{A}$. Entonces:
\begin{itemize}
\item $\displaystyle\sum P(A_n) < \infty \Longrightarrow P(\displaystyle \lim \sup A_n)=0$
\item Si $ A_n$ son independientes:
\begin{itemize}
\item $\displaystyle\sum P(A_n) < \infty \Longrightarrow P(\displaystyle \lim \sup A_n)=0$
\item $\displaystyle\sum P(A_n) = \infty \Longrightarrow P(\displaystyle \lim \sup A_n)=1$
\end{itemize}
\item $\displaystyle\sum P(A_n) < \infty \Longrightarrow P(\displaystyle \limsup A_n)=0$
\end{itemize}
\end{lemma}

\begin{proof}
  $$\limsup A_n=\bigcap_{n=1}^\infty \bigcup_{k=n}^{\infty}A_k=\bigcap_{n=1}^\infty B_n$$
  $$P(\limsup A_n)=P\lim B_n=\lim P\, B_n = \lim_n P\bigcup_{k=n}^\infty A_k \leq \lim_n \sum_{k=n}^\infty P\, A_k=0$$
  El recíproco no es cierto.
  

\end{proof}

\subsubsection{Criterio cero-uno de Borel}
\begin{criterion} 
  \ 
Sea $(\Omega,\mathcal{A},P) \, , \, A_n \in \mathcal{A}$. Entonces:

\begin{itemize}
\item Si $ A_n$ son independientes 
  \begin{itemize}
  \item $\displaystyle\sum P(A_n) < \infty \Longrightarrow P(\displaystyle \limsup A_n)=0$
  \item $\displaystyle\sum P(A_n) = \infty \Longrightarrow P(\displaystyle \limsup A_n)=1$
  \end{itemize}
\end{itemize}
\end{criterion}

\begin{proof}
  \
  
Usamos las sucesiones definidas anteriormente, para el caso 0 queda probado con el lema de Borel-Cantelli, para el caso 1, se tiene:
$$P(\displaystyle\limsup A_n)^c = P(\liminf A_n^c)=P(\bigcup_{n}\bigcap_{k=n}A_k^c)=P(\lim_{n}\bigcap_{k=n}A_k^c)=$$
$$\lim_n P(\bigcap_{k=n}A_k^c)=\lim_n \prod_{k=n}(1-P(A_k))\leq \lim_{n}\prod_{k=n} e^{-P(A_n)}=\lim_n e^{-\sum_k P(A_k)}=0$$


\end{proof}

De esta ley se extraen una serie de corolarios importantes:\\

\textbf{Corolario 1: }Si $A_n$ son independientes:
$$P(\lim \sup (A_n)) = 0 \Leftrightarrow \sum_{n\in\mathbb{N}} P(A_n) < \infty$$

\textbf{Corolario 2: }Sea ${X_n}$ una colección de variables aleatorias independientes, entonces;
$$ X_n \stackrel{c.s.}{\rightarrow} 0 \Leftrightarrow \sum_n P([X_n \geq \varepsilon]) < \infty$$

\textbf{Demostración}
  
$\sum P([|X_n| \geq \varepsilon]) = 0 
\ \Leftrightarrow \
P(\lim \sup ([|X_n| \geq \varepsilon])) = 0 
\ \Leftrightarrow \
\displaystyle P ( \bigcap_n \bigcup_{k=n} [|X_n| \geq \varepsilon]) = 0 
\ \Leftrightarrow \
P (\lim_n \bigcup_{k=n} [|X_n| \geq \varepsilon])
\ \Leftrightarrow \
\lim_{n \rightarrow \infty} P (\bigcup_{k=n} [|X_k| \geq \varepsilon]) = 0 
\ \Leftrightarrow \
X_n \stackrel{c.s.}{\rightarrow} 0$



\textbf{Corolario 3: }

${X_n}$ son independientes y $ X_n \stackrel{\mathbb{P}}{\rightarrow} 0 \Leftrightarrow \exists {X_{n_k}} \subset {X_n}$ tal que ${X_{n_k}} \stackrel{c.s.}{\rightarrow} 0 $ 

\

\textbf{Demostración}

$\exists X_{n_k}: P[|X_{n_k}| \geq \frac{1}{2^k}] \ \leq \ \frac{1}{2^k}$\\
$\displaystyle \sum_k P ([|X_{n_k}| \geq \frac{1}{2^k}]) \leq \sum_k \frac{1}{2^k} = 1 \Leftrightarrow X_{n_k} \stackrel{c.s.}{\rightarrow}0$




Pasamos ahora a ver otra ley importante en el estudio de la probabilidad, la cual se conoce como la ley 0-1 de Kolmogorov. Para ello, comenzamos definiendo el concepto de $\sigma$-campo cola.\\

Sean ${X_n}$ independientes, tenemos las siguientes relaciones entre los $\sigma$-campos asociados:\\

$\mathcal{B}(X_n, X_{n+1}) \supseteq \mathcal{B}(X_n)$\\
$\mathcal{B}(X_n, X_{n+1}, ...) \supseteq \mathcal{B}(X_{n+1}, ...)$ sucesión decreciente\\
$\mathcal{B}(X_n, X_{n+1}, ...), \mathcal{B}(X_{n+1}, X_{n+2}, ...), ... \downarrow \lim_n \sup \mathcal{B}(X_n) \equiv \mathcal{C}$\\

Dicho límite es un $\sigma$-campo, conocido como el $\sigma$-campo cola, que será equivalente a $\emptyset$ o $\Omega$.

\subsection{Ley 0-1 de Kolmogorov}
Sea $X_n$ una sucesión de variables aleatorias independientes. Entonces los eventos en la cola tienen o bien probabilidad 0 o 1 y las funciones cola son degeneradas.
\begin{proof}$\,$\\
$i)$\\
Por un lado, tenemos que: $\mathcal{C} \subseteq \mathcal{B}(X_{n+1}, ...)$\\
Además, $\mathcal{B} (X_1, ..., X_n)$ es independiente de $\mathcal{B}(X_{n+1}, X_{n+2},...)$, lo que implica que $\mathcal{C}, \mathcal{B}(X_1,...,X_n)$ son independientes. De aquí tenemos que $\mathcal{C}$ es independiente de $\mathcal{B}(X_1,...,X_n,X_{n+1},...)$, por lo que $\mathcal{C}$ es independiente consigo mismo, ya que está contenido en el $\sigma$-campo que hace medibles a todos. Una vez aquí, tenemos que:
$$ P(AA) = P(A)P(A) = P(A)\, \ \forall A \in \mathcal{C} \Rightarrow P(A) = \{0,1\}\,\ \forall A \in \mathcal{C}$$
$ii)$\\
Los límites de las sucesiones de variables aleatorias independientes son $\mathcal{C}$-medibles (también para limsup). Por tanto, es una variable aleatoria que degenera


\end{proof}

A raíz de esta ley, podemos extraer también un corolario.\\

\textbf{Corolario. }Sea $X_n$ una sucesión de variables aleatorias independientes, entonces $X_n$ converge c.s. o diverge c.s., igualmente para la serie $\sum_{n\in\mathbb{N}} X_n$


\subsection{Variables aleatorias independientes y espacios producto}
\section{Convergencia y estabilidad de sumas; centramiento y truncamiento}
Esta sección está dedicada a la investigación de las sumas de variables aleatorias independientes y especialmente, a sus propiedades en el límite. Primero de todo, consideraremos $S_n=\displaystyle\sum_{k=1}^n X_k$ una suma de variables aleatorias independientes $X_1, X_2,\ldots X_n$
\\\\
\textbf{Definición de estabilidad: }Si existe una sucesión $\{a_n\}$ y otra $\{b_n\} \rightarrow \infty$ de forma que $\dfrac{S_n}{b_n}-a_n\stackrel{\mathcal{P}}{\rightarrow}0$ o $\dfrac{S_n}{b_n}-a_n\stackrel{c.s.}{\rightarrow}0$, entonces decimos que la sucesión $S_n$ es estable en probabilidad o c.s. 
\\\\
A continuación presentamos una ley que ha marcado el desarrollo de toda la teoría de la probabilidad.


\subsection{Ley de Bernouilli de los grandes números} % Pag 244 del Loeve
Sean $X_1,\ldots X_n$ variables aleatorias independientes e idénticamente distribuidas tales que $P[X_n=1]=p$ y $P[X_n=0]=q=1-p$. Entonces se tiene que:
$$\dfrac{S_n}{n}-p \stackrel{\mathcal{P}}{\longrightarrow}0$$
El Teorema del Límite Central, que desarrollaremos más adelante, es una consecuencia directa del perfeccionamiento de esta ley por parte de Moivre y Laplace.

\subsection{Ley fuerte de Borel de los grandes números} % Pag 244 del Loeve
Con las hipótesis anteriores de la ley de Bernouilli de los grandes números, entonces: $$\dfrac{S_n}{n}-p \stackrel{c.s.}{\longrightarrow}0$$
\subsection{Teorema de aproximación de Weierstrass}
\begin{theorem}
Sea f(x) continua en $[0,1]$ se aproxima por el polinomio de Berstein $\Rightarrow \exists {B_n(x)}$ tal que $\lim B_n (x) = f(x)$ uniformemente $ \forall x \in [0,1]$.
\end{theorem}
\begin{proof}*
\end{proof}
\subsection{Centramiento y truncamiento}
Diremos que centramos $X$ a $c$ si reemplazamos $X$ por $X-c$. Si $X$ es integrable, entonces podemos centrarla a su esperanza $EX$ de forma que $X$ es reemplazada por $X-EX$. Es decir, diremos que una variable aleatoria está centrada a su media si y solo si su media existe y es igual a 0.
\\\\
Sea $X$ integrable, llamaremos varianza de $X$ al segundo momento de $X-EX$ (podría ser $\infty$) y lo denotaremos como $\sigma^2 X$. Por consiguiente:
$$\sigma^2 X = Var(X) = E(X-EX)^2=EX^2-(EX)^2$$
Para todo $c$ finito si centramos una variable en $c$ su varianza no cambia:
$$\sigma^2(X-c)=E(X-c-E(X-c))^2=E(X-EX)^2$$

Una vez visto el centramiento vamos con el \textbf{truncamiento}, útil cuando no existe la media.
\\\\
Diremos que truncamos $X$ a $c>0$ (finita) cuando reemplazamos $X$ por $X^c$ si $|X|<c$ o por $0$ si $|X|\geq c$, se define la variable $X^c$ como $X$ truncada a $c$. 
\\\\
Tenemos entonces que si $F$ es el diferencial de $X$, entonces existen todos los momentos de $X^c$ y son finitos; estos son:
$$EX^c=\int_{|X|<c}x\, dF\,\,\, , \,\,\, E(X^c)^2=\int_{|X|<c}x^2\, dF \,\,\, , \,\,\,  \ldots$$
Siempre podremos seleccionar un $c$ suficientemente grande para hacer que $P[X\neq X^c]=P[|X|\geq c]$ sea arbitrariamente pequeño. Además nosotros también podemos seleccionar el $c_j$ suficientemente grande para conseguir $P\bigcup [X_j\neq X_j^{c_j}]$ arbitrariamente pequeño, de tal forma que, dado $\epsilon > 0$, tenemos:
$$P\left( 
\bigcup [X_j\neq X_j^{c_j}] \right)
\ \leq \ 
\sum P[|X_j|\geq c_j] 
\ < \
\frac{\epsilon}{2^j}$$
Por consiguiente, dada una familia numerable de variables aleatorias, podemos hacer corresponder una familia de variables aleatorias que difiere de la primera en un evento de probabilidad pequeña. Además, si estamos interesados principalmente en propiedades en el límite, no necesitamos probabilidades pequeñas arbitrarias.
\\\\
\textbf{Definición. }Dadas dos sucesiones $X_n$ y $X_n^{\prime}$, se dice que son \textbf{cola-equivalentes} si difieren casi seguramente en un número finito de elementos.
En otras palabras, si $\forall \omega\in\Omega$ existe un número finito $n(\omega)$ tal que para $n\geq n(\omega)$ las dos sucesiones $X_n(\omega)$ y $X_n^{\prime}(\omega)$ son la misma.

\

$X_n$ y $X_n^{'}$ son cola-equivalentes si 
$P \left( \limsup [X_n \not = X_n^{'}] \right) = 0$.
\\\\
\textbf{Definición. } % Pag 245 Loeve
Se dice que $X_n$ y $X_n^{\prime}$ son  \textbf{equivalentes en la convergencia} si el conjunto de convergencia de ambos es el mismo c.s.



\begin{lemma} (de equivalencia)
  
Si las series $\sum_{n\in\mathbb{N}} P[X_n\neq X_n^{\prime}]$ convergen, entonces las sucesiones $X_n$ y $X_n^{\prime}$ son cola-equivalentes, por tanto, las series $\sum_{n\in\mathbb{N}} X_n$ y $\sum_{n\in\mathbb{N}} X_n^{\prime}$ son equivalentes en  convergencia y las series $\dfrac{S_n}{b_n}$ y $\dfrac{S_n^{\prime}}{b_n}$ con $b_n \rightarrow\infty$ convergen al mismo evento y al mismo límite, excluyendo el evento nulo.
\end{lemma}


% Sean $X_n$ variables aleatorias independientes truncadas en $E(X_n)$, entonces $$ Var\left( \sum_{k=1}^n \right) = \sum_{k=1}^n Var(X_k) $$

\

Si las variables aleatorias $X_n$ son independientes e integrables, entonces:
$$ \sigma^2 S_n=\sum_{k=1}^n \sigma^2 X_k \hspace{1.5cm} \text{(Igualdad de Bienaymé)}$$

\

\textbf{Demostración}

Sean $X_1$, $X_2$, ... integrables, como el centramiento no modifica la varianza, podemos asumir, cuando computamos varianzas, que esas variables aleatorias están centradas a la media. Entonces:
$$\sigma^2 S_n=ES_n^2=\sum_{k=1}^n EX_k^2+\sum_{j,k=1, j\not = k}^n EX_jX_k=\sum_{k=1}^n\sigma^2 X_k$$
como $X_j$ y $X_k$ son independientes, tenemos que:
$$EX_jX_k=EX_j\cdot EX_k=0$$




\

Entonces obtenemos la \textbf{desigualdad de Bienaymé-Tchebychev}:
$$\frac{\displaystyle\sum_{k=1}^n \sigma^2 X_k-\epsilon ^2}{c.s.sup(S_n-ES_n)^2} 
\ \leq \ 
P[|S_n-ES_n|\geq \epsilon]
\ \leq \ 
\frac{\sigma^2(S_n-ES_n)}{\epsilon^2} 
\ \leq \ 
\frac{1}{\epsilon^2}\sum_{k=1}^n\sigma^2 X_k$$
Esta desigualdad, aplicada a $(S_{n+k}-ES_{n+k})-(S_n-ES_n)$ y remplazando $\epsilon$ por $b_n$ y pasando al límite se tiene el siguiente corolario:
\\\\
\textbf{Corolario: }
\begin{itemize}
\item Si las series $\sum\sigma^2 X_n$ convergen, entonces las series $\sum (X_n-EX_n)$ convergen en probabilidad.
\item Si $\dfrac{1}{b_n^2}\displaystyle\sum_{k=1}^n \sigma^2 X_k\rightarrow 0$, entonces $\dfrac{S_n - ES_n}{b_n}\stackrel{\mathcal{P}}{\rightarrow}0$\\
\end{itemize}
\begin{proof}
% REVISAR DEMOSTRACIÓN
La demostración pasa por notar que:
$$P \left[\sum_{k}^{n+h}(X_k - EX_k)\geq\epsilon\right]
\ \leq \ \dfrac{\displaystyle\sum_{k}^{n+h}\sigma^2 X_k}{\epsilon^2}\longrightarrow 0$$
\end{proof}
Este último corolario es debido a Tchebichev (cuando $b_n=n$). En el caso de Bernouilli, donde $b_n=n$, $EX_n=p$ , $\sigma^2 X_n=pq$ se reduce a la ley de Bernouilli de los grandes números.

% DESIGUALDADES DE KOLMOGOROV (PAGINAS 31 A 34 DE MARIADEL)

\subsection{Ley débil de los grandes números}
Sea $\{X_n\}$ una sucesión infinita de variables aleatorias independientes que tienen el mismo valor esperado $\mu$ y varianza $\sigma^2$, entonces 
$$
 \overline{X}_n = \frac{X_1+\cdots+X_n}{n} 
 \stackrel{P}{\longrightarrow} \mu
$$

\subsection{Ley fuerte de los grandes números de Kolmogorov}
Sea ${X_n}$ una sucesión de variables aleatorias independientes e idénticamente distribuidas con una ley común $\mathcal{L}(X)$. Entonces:
$$E|X| < \infty
\Longleftrightarrow 
\dfrac{S_n}{n}=\dfrac{\sum_{k=1}^n X_k}{n}\stackrel{c.s}{\longrightarrow}EX$$

% Prueba página 251 del Loeve


% METODOS DE MONTECARLO

\chapter{El problema central del límite}
El problema del límite central en probabilidad es el problema de convergencia de leyes de secuencias de sumas de variables aleatorias.
\section{Preeliminares}
Consideramos $S_n$ el número de ocurrencias de un evento con probabilidad $p$ en $n$ intentos idénticos e independientes. Para evitar trivialidades asumimos que $pq\neq 0$, donde $q=1-p$. Si $X_k$ denota el indicador del evento en el $k$-ésimo intento, entonces $S_n=\displaystyle\sum_{k=1}^n X_k$ con $n=1,2,\ldots$ donde los sumandos son indicadores independientes e idénticamente distribuidos. Este es el caso de Bernouilli; como $EX_k=p$ y $EX_k^2=p$, entonces $\sigma^2 X_k=p-p^2=pq$ de donde se sigue que:
$$ES_n=\sum_{k=1}^nEX_k=np \,\,\, , \,\,\, \sigma^2 S_n=\sum_{k=1}^n \sigma^2 X_k=npq$$
El primer teorema del límite central publicado en 1713 dice que $\dfrac{S_n}{n}\stackrel{\mathcal{P}}{\longrightarrow}p$
\\\\
Puliendo este resultado, de Moivre obtuvo el segundo teorema de límites, que en la forma integral dada por Laplace, dice que:
$$P\left[ \frac{S_n -np}{\sqrt{npq}}< x \right]\rightarrow \frac{1}{\sqrt{2\pi}}\int_{-\infty}^{x}e^{-\frac{1}{2}y^2}dy\, ,\,\,\, -\infty\leq x\leq \infty$$
\\\\
El tercer teorema del límite central fue obtenido por Poisson, quien modificó el caso de Bernouilli asumiendo que la probabilidad $p=p_n$ depende del número de intentos de tal manera que $np_n\rightarrow\lambda>0$. Por consiguiente, escribiendo ahora $X_{nk}$ y $S_{nn}$ en lugar de $X_k$ y $S_n$, el caso de Poisson corresponde a una serie de sumas $S_{nn}=\displaystyle\sum_{k=1}^n X_{nk}$ donde, para cada n, los sumandos $X_{nk}$ son indicadores independientes e idénticamente distribuidos con $P[X_{nk}=1]=\dfrac{\lambda}{n}+o \left( \dfrac{1}{n}\right)$
\\\\
A partir de aquí Poisson probó que para $k=0,1,2,\ldots$
$$P[S_{nn}=k]\rightarrow \dfrac{\lambda^k}{k!}e^{-\lambda}$$
Así nacieron las tres leyes básicas de la teoría de la probabilidad:
\begin{enumerate}
\item La ley de degeneración $\mathcal{L}$ de una variable aleatoria degenera a 0 teniendo un punto de incremento solo en $x=0$
\item La ley normal $N(0,1)$ de una variable aleatoria normal con su función de densidad definida viene dada por:
$$F(x)=\frac{1}{\sqrt{2\pi}}\int_{-x}^{x}e^{-\frac{1}{2}y^2}dy$$
\item La ley de Poisson $\mathcal{P}(\lambda)$ de una variable aleatoria de Poisson con su función de densidad definida viene dada por:
$$F(x)=e^{-\lambda}\sum_{k=0}^x \frac{\lambda^k}{k!}$$
\end{enumerate}
Para comprobar que son válidas, debemos ver que el límite de las distribuciones es justamente la distribución del límite:
\begin{enumerate}
\item \textbf{Ley degenerada: }Tenemos que ver que: $\varphi_{X}(t)=E(e^{itX})=\int_{-\infty}^{\infty}=1$
$$\varphi_{\frac{S_n -ES_n}{n}}(t)=E(e^{it\frac{S_n-ES_n}{n}})=E(e^{it\sum^n X_k-EX_k\frac{1}{n}})=E(\prod_{k=1}^{n}e^{it(X_k-p)\frac{1}{n}})=$$
$$\prod_{k=1}^{n}E(e^{it\frac{X_k-p}{n}})=\prod_{k=1}^{n}(pe^{\frac{itq}{n}}+qe^{\frac{-itp}{n}})=(p(1+\frac{ip}{n}t+\ldots)+q(1-\frac{ip}{n}+\ldots))^n=$$
$$((p+q)+o(\frac{1}{n}))^n\rightarrow 1$$
\item \textbf{Ley normal: }Tenemos que ver que: $\varphi_X=E(e^{itX})=\frac{1}{\sqrt{2 \pi}}\int_{-\infty}^{x}e^{itX}e^{-\frac{x^2}{2}}dx=e^{-\frac{t^2}{2}}$
$$\varphi_{\frac{S_n-E(ES_n)}{\sqrt{npq}}}=E(e^{it\frac{\sum X_k-p}{\sqrt{npq}}})=\prod_{k=1}^{n} E(e^{it\frac{X_k-p}{\sqrt{npq}}})=\prod_{k=1}^{n}E(1+\frac{i(X_p-p)}{\sqrt{npq}}t+\ldots)=$$
$$\prod_{k=1}^{n}(1+\frac{i^2pq}{npq}\frac{t^2}{2}+\ldots)=(1+\frac{-t^2}{2n}+\ldots)^n=e^{-\frac{t^2}{2}}$$
\item \textbf{Ley de Poisson} $\varphi_X=E(e^{itX})=\displaystyle\sum_{x=0}^{\infty}e^{itX}e^{-\lambda}\frac{\lambda^x}{x!}=e^{-\lambda}\displaystyle\sum_{x=0}^{\infty}\frac{(e^{it}\lambda)}{x!}=e^{\lambda(e^{it}-1)}$
\end{enumerate}
Mientras que las dos primeras leyes juegan un papel central en el desarrollo de la teoría de la probabilidad, la de Poisson se ha mantenido ignorada durante mucho tiempo.
\section{Primer teorema del límite central (Problema clásico)}
En el caso de Bernouilli:
$$\mathcal{L}\left( \frac{S_n-ES_n}{n}  \right)\rightarrow \mathcal{L}(0) \,\, y \,\, \mathcal{L}\left( \frac{S_n-ES_n}{\sigma S_n}  \right)\rightarrow N(0,1)$$
y en el caso de Poisson:
$$\mathcal{L}(S_{nn})\rightarrow P(\lambda)$$
\begin{proof}

\end{proof}
Si consideramos variables aleatorias centradas en media, podemos enunciar el siguiente lema. Pero antes notamos $s_n=\sigma(S_n)$ y además tendremos en cuenta que las variables aleatorias están centradas en media, por lo que $EX_k=0$ y $ES_n=0$ a partir de ahora
\begin{lemma}
Si los sumandos son independientes, indénticamente distribuidos y centrados en media, entonces $\mathcal{L} \left( \dfrac{S_n}{n} \right)\rightarrow \mathcal{L}$ y $\mathcal{L} \left( \dfrac{S_n}{s_n} \right)\rightarrow N(0,1)$
\end{lemma}
\subsubsection{Teorema de composición y descomposición}
\begin{theorem}
Los tipos normales y degenerados son cerrados bajo composiciones y descomposiciones. Lo mismo también se cumple para cada familia de Leyes de Poisson $\mathcal{P}(\lambda ; a,b)$ con el mismo b
\end{theorem}
Una vez visto el lema anterior, enunciamos el siguiente teorema:

% FALTA LA PRIMERA EXTENSIÓN DEL TEOREMA

\begin{theorem} \textbf{Segunda extensión del problema clásico central del límite}\\

Consideramos $S_n=\displaystyle\sum_{k=1}^n X_k$ y $s_n=\sigma S_n$ donde los sumandos son variables aleatorias independientes centradas en media. Entonces:
\begin{enumerate}
\item Si $\dfrac{1}{n^{1+\delta}}\displaystyle\sum_{k=1}^n E|X_k|^{1+\delta}\rightarrow 0$ para $0\leq\delta\leq 1$. Entonces:
$$\mathcal{L}\left( \frac{S_n}{n} \right)\rightarrow \mathcal{L}(0)$$
\item Si $\dfrac{1}{s_n^{2+\delta}}\displaystyle\sum_{k=1}^n E|X_k|^{2+\delta}\rightarrow 0$ para un $\delta$ positivo, entonces:
$$\mathcal{L}\left( \frac{S_n}{s_n} \right)\rightarrow N(0,1)$$
\end{enumerate}
\end{theorem}
\begin{proof}
\end{proof}
Notar que si tomamos $\delta=1$, el teorema se convierte en el conocido Teorema de Tchevichev. Además también contiene al resultado de Markov:
\\\\
\textbf{Resultado de Markov:} Si $E|X_k|^{1+\delta}\leq c < \infty$, entonces $\mathcal{L}\left( \dfrac{S_n}{n} \right)\rightarrow \mathcal{L}(0)$, $\forall \delta >0$
\begin{lemma}
Si $S_{nn}=\displaystyle\sum_{k=1}^n X_{nk}$ donde los sumandos son variables aleatorias independientes (centradas en media), entonces:
\begin{enumerate}
\item Si $\dfrac{1}{n^2}\displaystyle\sum_{k=1}^n E|X_{nk}|^2\rightarrow 0$ entonces $\mathcal{L} \left( \dfrac{S_{nn}}{n} \right)\rightarrow \delta(0)$
\item Si $\dfrac{1}{s_{nn}^3}\displaystyle\sum_{k=1}^n E|X_{nk}|^3\rightarrow 0$ entonces $\mathcal{L} \left( \dfrac{S_{nn}}{s_{nn}} \right)\rightarrow N(0,1)$
\end{enumerate}
\end{lemma}
\subsubsection{Lema de equivalencia de leyes}
\begin{lemma}
Si $X_n - X_n ^{\prime}  \stackrel{\mathbb{P}}{\longrightarrow} 0$ o $P[X_n \neq X_n^{\prime}]\longrightarrow 0$ entonces las sucesiones $\mathcal{L}(X_n)$ y $\mathcal{L}(X_n^{\prime})$ son equivalentes.
\end{lemma}
Una vistos estos resultados, ya estamos en posición de poder dar una solución completa al problema clásico del límite central, pero antes enunciaremos:
\subsubsection{Criterio clásico de convergencia degenerada}
Se tiene que $\mathcal{L} \left( \dfrac{S_n}{n} \right)\rightarrow \mathcal{L}(0)$ si y solamente si:
\begin{enumerate}
\item $\sum\displaystyle\int_{|x|\geq n}dF_k \rightarrow 0$
\item $\dfrac{1}{n}\sum\displaystyle\int_{|x|<n} x dF_k \rightarrow 0$
\item $\dfrac{1}{n^2}\sum\left( \displaystyle\int_{|x|<n}x^2 dF_k - \left(  \displaystyle\int_{|x|<n}x\, dF_k \right)^2 \right)\longrightarrow 0$
\end{enumerate}
\subsubsection{Criterio clásico de convergencia normal}
Se tiene $\mathcal{L} \left( \dfrac{S_n}{s_n} \right)\rightarrow N(0,1)$ y $max_{k\leq n}\dfrac{\sigma_k}{s_n}\rightarrow 0$ si y solamente si, $\forall \epsilon >0$:
$$g_n(\epsilon)=\dfrac{1}{s_n^2}\sum\int_{|x|\geq \epsilon s_n}x^2\, dF_k\longrightarrow 0$$
\begin{proof}
.
\end{proof}

\newpage

\part{Ejercicios}
%%%%%%%%%%%%%%%%%%%%%%%%%%%%%%%%%%%%%%%%%%%%%%%%%%%%%%%%%%%%%%%%
%
% Documento con ejercicios resueltos de la
% asignatura Análisis Matemático I.
% Doble Grado de Informática y Matemáticas.
% Universidad de Granada.
% Curso 2016/17.
%
% Licencia:
% CC BY 4.0 (https://creativecommons.org/licenses/by/4.0/)
%
%%%%%%%%%%%%%%%%%%%%%%%%%%%%%%%%%%%%%%%%%%%%%%%%%%%%%%%%%%%%%%%

%--------------------------------------------
%   Topología de un espacio métrico
%--------------------------------------------
\section{Topología de un espacio métrico.}

\begin{ejer}
Dado el conjunto $A = \{ (x,y)\in \mathbb{R}^2: 0 < x \le 1 \}$, ¿es abierto?

\begin{proof}
Tenemos que comprobar si $A$ es abierto, es decir, si es cierto que\\
$\forall a \in A \ \ \exists s>0\ tal\ que\ B(a,s)\subset A$. Para ello, fijo $y \in \mathbb{R}$ y escojo $x_o = (1,y) \in A$. Además, tomo $s>0$ cualquiera. Veamos que hay puntos $z \in B(x_o, s)$ que no pertenecen a $A$.

Sea $\displaystyle z = (1 + \frac{s}{2}, y)$. Entonces, se tiene que $\displaystyle d(z,x_o)= \sqrt{\left(\frac{s}{2}\right)^2 + 0} = \frac{s}{2} < s$, y por tanto $z \in B(x_o,s)$.

Claramente $z\notin A$, pues $\displaystyle 1 + \frac{s}{2} > 1$. Así, concluimos que $z$ es un punto de $B(x_o,s)$ que no pertenece a $A$, por lo que $\displaystyle B(x_o,s) \not \subseteq A$, y $A$ no es abierto. \qedhere

%\textbf{@ Antonio Coín.}
\end{proof}

\end{ejer}

\section{Conexión.}
\begin{ejer} Dado el conjunto \emph{peine} de la siguiente forma, probar que es conexo:
\[
	P = (\{ 0\} \times (0,1]) \cup ((0,1] \times \{0\}) \cup (\{\frac{1}{n} : n \in \mathbb{N}\} \times (0,1))
\]

\begin{center}
\includegraphics[scale=0.4]{./img/peine.png}
%\includegraphics[scale=0.4]{./Análisis Matemático I/imagenes/peine.png}
%% Descomentando la línea anterior se puede generar el pdf con el rakefile
%% Dejo sin comentar la primera porque normalmente trabajamos con el .tex desde la carpeta de Análisis
\end{center}

\begin{proof}
Para la demostración vamos a definir el siguiente conjunto:
\[
	S = ((0,1] \times \{0\}) \cup (\{\frac{1}{n} : n \in \mathbb{N}\} \times (0,1))
\]
Es decir, el mismo conjunto P quitandole el eje y. Siendo $\bar{S} = P \cup (0,0)$.
De esta forma tenemos:
\[
S \subset P \subset \bar{S}
\]

Para demostrar este ejercicio utilizaremos la siguiente proposición:\\

\begin{nprop}
Sea A un conjunto conexo, y B tal que $A \subset B \subset \bar{A} \implies B$ es conexo.
\end{nprop}

\begin{proof}
 Para la demostración probaremos su contrarecíproco, consistente en: B no es conexo $\implies$ A no es conexo.

B no es conexo $\implies \exists O,O'$ abiertos $: \begin{cases}
B \subset O\cup O'\\
O \cap B \cap O' = \emptyset\\
O \cap B \neq \emptyset \neq O' \cap B\\
\end{cases}$

 Debemos probar que los 3 resultados se cumplen para A.

 1. $B \subset O \cup O' \implies A \subset B \subset O \cup O'$\\
 2. $B \cap (O\cap O') = \emptyset \implies \nexists x \in B$ tq $x \in O \cap O'$. \\
 Entonces $\forall x \in A \implies x \in B \implies x \not\in O \cap O'$\\
 3. $O \cap B \neq \emptyset \implies O\cap \bar{A} \neq \emptyset \implies O\cap (A\cup fr(A) \neq \emptyset \implies (O\cap A)\cup (O\cap fr(A)) \neq \emptyset$
 En este punto tenemos 3 posibilidades:\\
  - El primer elemento de la union es distinto del vacio, en cuyo caso hemos acabado.\\
  - Ambos elementos son distintos del vacio, en cuyo caso tambien queda demostrado.\\
  - El segundo elemento de la unión es distinto del vacio. $\exists x \in O\cap fr(A) \implies \exists \epsilon > 0 : B(x,\epsilon) \in O$. Por estar en la frontera: $B(x,\epsilon)\cap A \neq \emptyset \implies \exists y \in B(x,\epsilon)\cap A \subset A \cap O$.\\\\

  Nota: La idea de esta demostración es que: dado un conjunto conexo, su unión con elementos que se encuentran cerca de el, sigue siendo conexa.

\end{proof}

Teniendo esto demostrado, haciendo uso de la proposición y de que S es conexo (por ser arco-conexo). Queda demostrado que P es conexo.

\end{proof}

\end{ejer}


\begin{ejer}
	Estudiar los máximos y mínimos relativos de la función $f: R^2 \to R$ dada por:
	\begin{enumerate}

	\item $f(x,y)=x^2+y^2$. Tenemos que estudiar los puntos críticos de f. Para ello, estudiamos:
	\[
	\begin{rcases}
	\frac{\partial f}{\partial x}(x,y) = 0\\
	\frac{\partial f}{\partial y}(x,y) = 0
\end{rcases} \implies \begin{rcases}
	2x= 0\\
	2y = 0\\
\end{rcases}
	\]
	Ahí podemos obtener los puntos críticos, que como vemos es el (0,0). Ahora, tenemos que ver si es un máximo o un mínimo usando la matriz Hessiana, que se calcula derivando dos veces por cada variable.[REVISAR CÓMO SE OBTENDRÍA LA HESSIANA]

	Calculamos la matriz Hessiana y la evaluamos en el (0,0), el punto obtenido:
	\[
	Hf(x,y) = \begin{pmatrix} 2 & 0 \\ 0 & 2 \end{pmatrix}
	\]
	Por tanto en el (0,0) también es esa matriz y queda ver que es definida positiva, pero eso es trivial por los determinante encajados. Por tanto, $f$ tiene un mínimo local en 0.

	\item $f(x,y) = -x^2-y^2$. Este apartado es igual, pero la matriz tendrá -2 en vez de 2 y por ello tendrá un máximo.

	\item $f(x,y) = x^2-y^2$. Estudiamos puntos críticos:
	\[
	\begin{rcases}
	\frac{\partial f}{\partial x}(x,y) = 0\\
	\frac{\partial f}{\partial y}(x,y) = 0
\end{rcases} \implies \begin{rcases}
	2x= 0\\
	-2y = 0\\
\end{rcases}
	\]

	Por tanto, su matriz hessiana es:
		\[
	Hf(x,y) = \begin{pmatrix} 2 & 0 \\ 0 & -2 \end{pmatrix}
	\]

	Y esta matriz no es ni semidefinida positiva ni semidefinida negativa. Por tanto, en este caso el (0,0) no es ni máximo ni mínimo relativo.
	Como no es ni máximo ni crítico pero su derivada es cero, se denomina punto de silla.

	\item $f(x,y)=x^2-xy+y^2$ Estudiamos puntos críticos:
	\[
	\begin{rcases}
	\frac{\partial f}{\partial x}(x,y) = 0\\
	\frac{\partial f}{\partial y}(x,y) = 0
\end{rcases} \implies \begin{rcases}
	2x -y= 0\\
	2y-x = 0\\
\end{rcases}
	\]

	Resolvemos el sistema, y nos queda la solución (0,0) también como solución única.

	Por tanto, su matriz hessiana es:
		\[
	Hf(x,y) = \begin{pmatrix} 2 & -1 \\ -1 & 2 \end{pmatrix}
	\]

	Como los determinantes encajados son positivos, la matriz hessiana es definida positiva por tanto (0,0) es un mínimo


\end{enumerate}
\end{ejer}

\begin{ejer}[5 relación derivabilidad]
	Sea $\Omega \subset \R^n$ abierto y $f:\Omega \to \R^n$ y $a\in \Omega$ con $Df(a) = 0$ ( es decir, $\triangledown f(a) = 0$) con $def(Hess f(a)) \ne 0$.

	Probar que $\exists U$ abierto en $\Omega$ tal que $a \in U$ tal que:
	\[
	\begin{rcases}
	Df(x) = 0\\
	x \in U
\end{rcases} \iff x = a
	\]
\end{ejer}
\begin{proof}[Solución]
	Calculamos su hessiana:
	\[
	Hf(a) = \begin{pmatrix} \frac{\partial^2f(a)}{\partial^2x_1^2}& \cdots & \frac{\partial^2f(a)}{\partial x_1 \partial x_n} \\ \cdots & & \cdots \\ \frac{\partial^2f(a)}{\partial x_n \partial x_1}& \cdots & \frac{\partial^2f(a)}{\partial^2x_n^2} \end{pmatrix} = Jg(a) = J(\triangledown f)(a)
	\]
Y tenemos que $g:\Omega \to \R^n$, $g(x) = \triangledown f(x) = ( \frac{\partial f(x)}{\partial x_1}, \cdots , \frac{\partial f(x)}{\partial x_n})$.

Ahora, como el determinante de la Hessiana de f en a es distinto de cero, por el Teorema de la función inversa aplicado a $g =  \triangledown f$, entonces $\exists U$ abierto en $\Omega$, $a\in U$ y también $\exists V$ abierto en $\R^n$ con $g(a) =  \triangledown f(a) \in V$
tales que:

$g = \triangledown f: U \to V$ biyectiva que lleva $a \mapsto g(a) =  \triangledown f(a) = 0$


\end{proof}
\begin{ejer}
	Sea $f:\R^2 \to \R^2$ dada por $f(x,y) = (e^x cos(y),e^x sen(y)) = e^x(cosy,seny) \in \partial B(0,e^x)$. Probar que $f$ es localmente invertible en todo punto de $\R^2$
\end{ejer}
\begin{proof}[Solución]
	Primero, es fácil ver que $f$ no es inyectiva en todo $\R^2,$pues basta ver que $f(x,y+2\pi) = f(x,y)$

	Lo único que tendríamos que hacer es ver que el determinante de $Jf (x_0,y_0)$.
	\[
det \begin{pmatrix} e^x cosy & -e^xseny \\ e^x seny & e^x cos y \end{pmatrix}  = e^x \ne 0
	\]
	Y como no es 0, la función es localmente invertible en el punto $x_0,y_0$ por el Teorema de la función inversa.

\end{proof}
\begin{ejer}[6 de la relación]

\[
\begin{rcases}
	x^2 - y^2 -u^3 +v^4 = -4\\
	2xy + y^2 -2u^2 -3v^2 = -8
\end{rcases} = \begin{rcases}
	x^2 - y^2 -u^3 +v^4  + 4= 0\\
	2xy + y^2 -2u^2 -3v^2  + 8=  0
\end{rcases}
\]
	Podemos considerarla como la función $F: \R^4 \to \R^2$ dada por\\
	 $F(x,y,u,v) = (x^2 - y^2 -u^3 +v^4  + 4,2xy + y^2 -2u^2 -3v^2  + 8)$
	 Y podemos ver la ecuación como $F(x,y,u,v) = (0,0)$.
	 Tendremos que usar dos variables como parámetros, pues necesitaríamos 4 variables para resolver el sistema.
	 Hacemos entonces la jacobiana de esta función:

	 \[
	 det\begin{pmatrix}
 2x & -2y \\
 2y & 2x+2y
\end{pmatrix}
	 \]
	 Y habría que evaluar el determinante en el punto que nosotros habíamos tomado, que es una solución particular que sería : $x=2, y = 1,u = 2,v = 1$, por lo que sería ver el determinante de:
	 \[
	 \begin{pmatrix}
 4 & 2 \\
 -2 & 0
\end{pmatrix}
	 \]
, que es distinto de 0, por lo que por el teorema de la función implícita, existen U entorno abierto de (2,-1) en $\R^2$ y existe V entorno abierto de (2,1) en $\R^2$ tales que:
\begin{enumerate}
	\item UxV$\subset A = \R^2$
	\item $\forall(u,v) \in V \ \ \exists ! (x,y) \in U : F(x,y,u,v) = (0,0)$
\end{enumerate}
\end{ejer}



\begin{ejer}
	Sea $f:\R^2 \to \R$
	$f(x,y) = 2x^2 -y^2 +4x+2y+1$.

	\textbf{Procedimiento}: Para estudiar los extremos relativos hacemos puntos criticos y eso nos lleva a un conjunto de puntos que tendremos que estudiar. Luego,si estudiamos el máximo o el mínimo absoluto en un compacto, tendremos probar que el máximo y el mínimo se alcanzan (T.W.) y sabiendo que son máximo o mínimos absolutos, si el punto estuviera en el interior sería un máximo o mínimo relativo luego tenemos que estudiar los puntos de los bordes estudiando las funciones que están definidas en un abierto(IMPORTANTE). Cuando tengamos los puntos que puedan ser puntos críticos, calculamos su valor por la función y afirmamos cuál es el máximo y cuál es el mínimo.\\

\textbf{Apartado a:}\\
	 Estudiar extremos relativos y puntos de silla de f.

	Los puntos criticos de f son los puntos en los que el gradiente de f vale 0 ( $\triangledown f (x,y) = (0,0)0)$. Calculamos derivadas:
	\[
	\frac{\partial f}{\partial x}(x,y) = 4 x+ 4
	\]\[
	\frac{\partial f}{\partial y}(x,y) = -2y +2
	\]
	Por lo que tenemos que calcular las soluciones del sistema:
	\[
	\begin{rcases}
	4x+4 = 0\\
	-2y+2 = 0
\end{rcases}
	\]
	Cuya solución es el punto (-1,1).
	Estudiamos la matriz hessiana en ese punto para ver qué es.

	\[
	Hesf(x,y) = \begin{pmatrix}
 \frac{\partial^2f}{\partial x^2}(x,y) &  \frac{\partial^2f}{\partial x \partial y}(x,y) \\
  \frac{\partial^2f}{\partial x \partial y}(x,y) &  \frac{\partial^2f}{\partial y^2}(x,y)
\end{pmatrix}  =  \begin{pmatrix}
  4& 0 \\
 0 & -2
\end{pmatrix}
	\]
	Y el determinante de esa es -8 y además el primer determinante encajado es positivo, luego esa matriz no es ni definida positiva ni definida negativa luego tenemos que en (-1,1) hay un punto de silla (esta es su definición).

	\textbf{Apartado b:}

	Si llamamos T= triángulo cerrado de vértices (-3,0),(0,0) y (0,3) que es cerrado y acotado. Este T es compacto por ser cerrado y acotado por tanto, por el teorema de Weierstrass f alcanza un máximo valor y un mínimo en puntos de T.
	Esto es un problema de extremos condicionados.

	Si estuvieran los puntos críticos que hemos encontrado antes dentro del triángulo, entonces serían extremos relativos y hemos visto antes que no tiene. Luego el máximo y mínimo absoluto de f en T no puede ser un punto interior a T. Hay que estudiarlos en la frontera de T.

	La frontera de T está constituida por los 3 segmentos que unen los 3 puntos del triángulo y los puntos, es decir, "6 partes". Esto es:
	\[
	\partial T = A_1 \cup A_2 \cup A_3 \cup (0,3) \cup (-3,0) \cup (0,0)
	\]
	Siendo:\\ $A_1 = \{ (0,y) : 0 < y < 3\}  = \{(x,y) \in (-3,0)x\R : x = 0\}$,\\
	 $A_2 =\{(x,0) :  -3 < x < 0\} = \{(x,y) \in (-3,0)x \R : y = 0\}$  \\
	 $A_3 =\{(x,x+3) :  -3< x < 0\} = \{(x,y)\in (-3,0)x\R :  y -x-3 = 0\}$.\\

	Ahora, al descomponerlo así podemos usar en cada fragmento el teorema de lso multiplicadores de Lagrange.

	Empezamos por $A_3$ . Definimos $g(x,y) = y-x-3$ y sabemos que los posibles extremos condicionados serían extremos de la Lagrangiana es decir, tienen que ser puntos críticos de la función: $f+\lambda g$ dependiendo de $\lambda$.
	Esta función es:
	\[
	F(x,y,\lambda) = f(x,y) + \lambda g(x,y)
	\]
	E igualamos de nuevo el gradiente de esta función a 0:
	\[
	\frac{\partial F}{\partial x} =  \frac{\partial F}{\partial y} =\frac{\partial F}{\partial \lambda} = 0
	\]
	Obteniendo el sistema:
	\[
	\begin{rcases}
	4x +4 - \lambda = 0\\
	-2y +2 + \lambda = 0\\
	 y-x-3 = 0
\end{rcases}\implies  \begin{rcases}
	x = 0\\
	y = 3 \\
\end{rcases}
	\]
	Pero el punto (0,3) no está en $A_3$, es un punto aparte luego no hay ningún extremo condicionado a la restricción de pertenecer a $A_3$. Si existiera un extremo relativo en ese segmento también sería extremo relativo sobre ese segmento, pero como hemos visto que no tiene, no puede haber extremo relativo.

	Hacemos lo mismo con $A_2$. En este caso $g(x,y) = y$ Nos queda el sistema:
	\[
	\begin{rcases}
	4x +4  = 0\\
	-2y +2 + \lambda = 0\\
	 y = 0
\end{rcases}\implies \begin{rcases}
	x=-1\\
	y = 0\\
	\lambda = -2
\end{rcases}
	\]
	Tenemos el punt (-1,0) que sí pertenece a $A_2$, por tanto este punto puede ser extremo relativo condicionado por la condición $y = 0$

	Repetimos el proceso para $A_1$. En este caso, $g(x,y) = x$. Volvemos a hacer el sistema y nos queda:
	\[
	\begin{rcases}
	4x +4 +\lambda = 0\\
	-2y +2  = 0\\
	 x = 0
\end{rcases}\implies \begin{rcases}
	x=0\\
	y = 1\\
	\lambda = -4
\end{rcases}
	\]
	Tenemos el punto(0,1) que sí está en este segmento.

	 Todo esto sirve para afirmar que el máximo y el mínimo de f en T es uno de estos puntos:
	\[
	(0,1),(-1,0),(0,0),(0,3),(-3,0)
	\]
	Ahora, sabiendo esto calculamos la f en cada uno de estos puntos y vemos cuáles son los extremos y hemos acabado.


\end{ejer}


\begin{ejer}
	Tomemos la gráfica  $G\{(x,1/x): x > 0\}$ (gráfica de la función $y= 1/x$) .Probar que existe un punto en G tal que
\[
d((0,0),(x,1/x) = Min\{d((0,0),(y,1/y) : y > 0\}
\]

Para ello, debemos definir
\[
f(y) = d((0,0),(x,1/x)) \ \ continua
\]
Si nos enfrentaramos a un ejercicio donde el conjunto no es compacto, un "truco" a seguir es coger un punto del conjunto al azar, y trazar la bola cerrada de distancia la del (0,0) a ese punto. Al tomar la interseccion entre el conjunto y esa bola cerrada, tenemos un conjunto cerrado, que nos aseguramos que contiene el mínimo.
\end{ejer}


\begin{ejer}[26 de la relacion de derivabilidad]
	Sea $\Omega \subset \R^n$ abierto y acotado . Sea $u \in \varphi(\bar{\Omega} \cup \varphi^2(\Omega)$.

	Apartado a:

	$\triangle u(x):= \sum_i^n \frac{\partial^u}{\partial x_i^2}(x) < 0 \ \forall x \in \Omega$ es la traza. Entonces, u no puede alcanzar su mínimo absoluto en $\bar{\Omega}$ en un punto de $\Omega$.

Nota:
	$\Omega \subset \R^n$ abierto y acotado $\implies \begin{rcases}
	cerrado \ + \ acotado = \ compacto \\
	u \in \mathcal{C}(\bar{\Omega})
\end{rcases}\implies $ u alcanza su maximo y minimo en $\bar{\Omega}$


\begin{proof}[Solución apartado a:] \textbf{Este ejercicio es trivial.}\\
	Supongamos que fuese falso, es decir:supongamos que el mínimo absoluto de $u$ en $\bar{\Omega}$ se alcanza en $a\in \Omega$ : $u(a) = min\{u(x): x \in \bar{\Omega}\}$, esto implica que $\exists r > 0 : B(a,r) \subset \Omega $ y $u(x) \geq u(a) \forall x \in B(a,r) \implies a$ es un mínimo relativo de u, por lo que:
	\[
	\frac{\partial u}{\partial x_i}(a) = 0 \quad \forall i = 1,...,N
	\]
	\[
	\frac{\partial^2 u}{\partial x_i^2}(a) \geq 0 \implies \sum_i^n \frac{\partial^2u}{\partial x_i^2} \geq 0
	\]
	Y esto es una contradicción.
\end{proof}

Apartado b:

Si el $\triangle u(x) \leq 0 \quad \forall x \in \Omega \implies (min)_{x\in \bar{\Omega}} \ \ u(x) = (min)_{x \in \partial \Omega}  \ \ u(x)$.


\end{ejer}

\begin{ejer}[14 de la relación de derivabilidad]\hfill\\
	Sea $f: \R^2 \to \R^2$ con $f(x,y) = (y+cosx,x+e^y)$\\
	Sea $g: \R^2 \to \R$ con $g(t,u) = t+u$
	\begin{nlist}
	\item Calcular la matriz jacobiana de $F=gof$
	\item Calcular la derivada direccional de $F$ en el punto $(0,0)$ a lo largo de la dirección $(2,-1)$
	\item Calcular el valor máximo de la derivada direccional de $F$ en $(0,0)$
	\item $F(x,y) = 2$ define a y como una función implícita diferenciable $\mathcal{C}^\infty$ en un entorno de (0,0).
\end{nlist}
\end{ejer}
\begin{proof}[Solución]
	\begin{nlist}
	\item $F=g(f(x,y)) = g(y+cosx,x+e^y) = y+cosx+x+e^y$.

	Lo haremos usando la regla de la cadena. $DF(x,y) = Jg(f(x,y))Jf(x,y)$
	\[
	Jf(x,y)= \begin{pmatrix}
 -senx & 1 \\
 1 & e^y
\end{pmatrix} = Jg(y+cosx,x+e^y)
	\]
	\[
	Jg(f(x,y))=Jg(y+cosx,x+e^y) = (1,1)
	\]
	\[
	Jg(f(x,y))Jf(x,y) = (1 \ \ \ 1)\begin{pmatrix}
 -senx & 1 \\
 1 & e^y
\end{pmatrix} = (1- senx \ \ \ 1+e^y)
	\]

	\item Vamos a calcular primero el vector unitario de (2,-1). $|(2,-1)| = 5 \implies (2/\sqrt5, -1/ \sqrt5)$. Ahora, tenemos
	\[
	D_{ (2/\sqrt5, -1/ \sqrt5)}F(0,0) = JF(0,0) \begin{pmatrix}
 2/\sqrt 5  \\
 -1/\sqrt5
\end{pmatrix} = 0
	\]

	\item \[
	D_vF(0,0) = JF(0,0) v =  \triangledown F(0,0) v \leq \sqrt 5 \ \ \ \forall |v| = 1	\]
	Es una cota superior por el lema de cauchy schwarz.
	Ahora, \[
	v = \frac{\triangledown F(0,0)}{|\triangledown F(0,0)|} \implies D_v F(0,0) = \triangledown F(0,0)  \frac{\triangledown F(0,0)}{|\triangledown F(0,0)|} = |\triangledown F(0,0)| = \sqrt 5
	\]

	Podemos obtener como un resultado o propiedad:

	\begin{nprop}
	Si $\triangledown F(a,b) \ne (0,0)$ entonces:
	\[
	D_{ \frac{-\triangledown F(a,b)}{|\triangledown F(a,b)|}}F(a,b) =  |\triangledown F(a,b)| \leq D_v F(a,b) \leq |\triangledown F(a,b)| = D_{ \frac{\triangledown F(a,b)}{|\triangledown F(a,b)|}}F(a,b) \ \ \ \forall |v| = 1
	\]
\end{nprop}

Otra forma de haber hecho el ejercicio era mediante extremos condicionados:
\[
(1 \ \ 2) \begin{pmatrix}
 v_1 \\
 v_2
\end{pmatrix}
\]

\item Lo haremos pensando en la función $H(x,y) =F(x,y)-2 = 0$, pues nos ayudará para usar el teorema de la función implicita. Ahora, la función H es de clase infinito pues es composición de dos funciones que son de clase infinito, por lo que podemos usar el teorema. Tenemos que comprobar que $H(0,0) = 0$, que se comprueba fácilmente.
\[
JH(x,y) = JF(x,y) = \begin{pmatrix}
 1-senx & e^y+1 \\

\end{pmatrix}
\]
\[
JH(0,0)= \begin{pmatrix}
 1 & 2 \\
\end{pmatrix}
\]
Ahora, tenemos que aclarar cuál de los menores vamos a coger de la matriz (1   2), pues pueden ser la x o la y las funciones implícitas. Para ver cual cogeremos, usar una regla pnemotécnica.

Calcular $y'(0)$ e $y''(0)$ con y la función implícita una vez que la hemos calculado.
\end{nlist}
\end{proof}

\begin{ejer}[14 de la relación de derivabilidad]
Sea $F:\mathbb{R}^2 \longrightarrow \mathbb{R}$ la función dada por $F(x,y) = y + cosx + x + e^y \quad \forall(x,y)\in \mathbb{R}^2$.

\textbf{Apartado d)}. Probar que $F(x,y)=2$ define a $y$ como función implícita diferenciable $C^{\infty}$ en un entorno de $(0,0)$. Además, calcular $y'(0)$ e $y''(0)$.

\begin{proof}[Solución]
Definimos $H:\mathbb{R}^2 \longrightarrow \mathbb{R}$ dada por $H(y,x)=F(y,x)-2$.

\textbf{Observación 1.} $H$ es $C^{\infty}$, por ser suma de funciones $C^{\infty}$.\\
\textbf{Observación 2.} $H(0,0) = 0 + cos0 + 0 + e^0 - 2 = 0$. Esto es condición indispensable para poder aplicar el \textit{teorema de la función implícita}.\\
\textbf{Observación 3.} Vemos el conjunto de definición de $H$ como $\mathbb{R}^2 = \mathbb{R} \times \mathbb{R}$, donde la primera componente es "\textit{el parámetro"} ($y$), y la segunda componente es \textit{"la variable"} ($x$).

El siguiente paso es comprobar que $\displaystyle \det \begin{pmatrix}
	\pder[H]{y}(0,0)\\
\end{pmatrix} \ne 0$:

$$\pder[H]{y}(y,x) = 1 + e^y;\quad \det \begin{pmatrix}
	\pder[H]{y}(0,0)\\
\end{pmatrix} = \det \left( 1 + e^0 \right) = \det (2) = 2 \ne 0.$$

Ahora, estamos en las condiciones de aplicar el \textbf{teorema de la función implícita} a $H$:
$\exists U$ entorno abierto de $0$ en $\mathbb{R}$, $\exists V$ entorno abierto de $0$ en $\mathbb{R}$, tal que:

\begin{enumerate}
\item $U\times V \subseteq \mathbb{R}^2$
\item $\forall x \in V \ \exists ! \ y \in U\ tal\ que \ H(y,x) = 0$
\end{enumerate}

Es decir, podemos definir la función $y:V \longrightarrow U$ que a cada $x\in V$ le hace corresponder un $y(x)\in U$, de tal forma que $H(y(x),x) = 0$. Además, el teorema garantiza que $y$ es $C^{\infty}$ por serlo $H$. Por tanto, hemos probado que $F(x,y)=2$ define a $y$ como función implícita $C^{\infty}$ en un entorno de $(0,0)$.

\underline{Calculamos ahora $y'(0)$}. Sabemos que $y(0)=0$, pues $H(0,0) = 0$. También sabemos que, dado $x\in V$, la pareja $(y(x),x)$ es solución de la ecuación $H(y,x)=0$. Por tanto, se tiene que:
\vspace{0.5em}
$$y(x) + cosx + x + e^{y(x)} - 2 = 0$$

Derivamos ahora la función con respecto a "la variable" $x$, obteniendo:
\vspace{0.5em}
$$y'(x) - senx + 1 + e^{y(x)}\cdot y'(x) = 0$$

Por último, evaluamos en el punto $x=0$, que sabemos que es solución:
$$y'(0) - sen0 + 1 + e^{y(0)}\cdot y'(0) = 0 \Rightarrow y'(0) + 1 + y'(0) = 0 \Rightarrow \bm{y'(0) = \frac{-1}{2}}$$

\underline{Calculamos ahora $y''(0)$}. Sabiendo que $y'(0) = -1/2$, volvemos a derivar la expresión anterior, obteniendo:
\vspace{0.5em}
$$y''(x) - cosx + e^{y(x)}\cdot y''(x) + e^{y(x)}\cdot y'(x)= 0$$

Por último, evaluamos en el punto $x=0$, que sabemos que es solución:
$$y''(0) - cos0 + e^{y(0)}\left( y''(0) + y'(0) \right) = 0 \Rightarrow y''(0) - 1 + y''(0) - \frac{1}{2} = 0 \Rightarrow \bm{y''(0) = \frac{3}{4}}$$
\end{proof}
\end{ejer}

\begin{ejer}[Función lipschitziana respecto de una variable] \hfill

En orden a probar el \textbf{teorema de Ascoli-Arzelá} (ó \textit{teorema de Picard-Lindelöf}) necesitábamos verificar que una cierta función $f$ cumplía la condición de Lipschitz respecto de su segunda variable, bajo la hipótesis de que $f$ era $C^1$. Sin embargo, basta que $f$ cumpla una condición menos restrictiva, como veremos.

\textbf{Proposición.} Sea $f:\mathbb{R}^2 \longrightarrow \mathbb{R}$ una función continua, derivable respecto de su segunda variable, y verificando que $\pder[f]{x}(s,x)$ es continua. Sean $a,b \in \mathbb{R}$ con $a < b$, y consideremos dos funciones $y,z \in C([a,b],\mathbb{R})$. Entonces, se tiene que: $$\exists K > 0 : \ \ |f(s,y(s)) - f(s,z(s)) | \le K|y(s)-z(s)| \ \ \forall s \in [a,b]$$

\begin{proof} \hfill

Para probar el resultado, podemos considerar sin pérdida de generalidad que $f$ está definida en el conjunto $[a,b] \times \mathbb{R}$. Además, como $y,z$ son funciones continuas definidas en un compacto, sabemos que sus imágenes son compactas. Así, podemos definir $A = y([a,b]) \cup z([a,b])$, y afirmamos que $A$ es compacto, pues es unión finita de compactos. Por tanto, existirán números reales $c$ y $d$ tales que $c = \text{mín} A$ y $d = m\acute{a}x A$, con $c \le d$. Entonces, podemos restringir aún más el conjunto de definición de $f$, y consideramos que está definida en el compacto $D = [a,b]\times [c,d]$.

Ahora, como $f$ es $C^1$ respecto de su segunda variable, tenemos que la función $\displaystyle\pder[f]{x}: D \longrightarrow \mathbb{R}$ es una función continua definida en un compacto. Por tanto, su imagen está acotada, y esto nos permite definir: $$M := sup \left\{ \left| \pder[f]{x}(s,x) \right|\colon (s,x) \in D\right\} \ge 0$$

Fijamos $s \in [a,b]$, y consideramos la función $f_s(x):[c,d] \longrightarrow \mathbb{R}$ dada por $f_s(x) = f(s,x)$.\\ Es inmediato comprobar que $f_s$ es $C^1$, pues $f_s'(x) = \pder[f]{x}(s,x)$. Observamos también que si $y(s)=z(s)$, la proposición se cumple trivialmente para cualquier valor de $K$, por lo que podemos suponer $y(s) < z(s)$ (en otro caso, intercambiamos sus papeles). Entonces, por el \textbf{teorema del valor medio} aplicado a $f_s$, tenemos que $\exists c \in [y(s), z(s)]$ verificando: $$|f_s(y(s)) - f_s(z(s))| = |f'_s(c)|\cdot|y(s)-z(s)|$$

De esta igualdad concluímos que: $$|f(s,y(s)) - f(s,z(s))| = \left|\pder[f]{x}(s,c)\right|\cdot|y(s)-z(s)| \le M\cdot|y(s)-z(s)|$$

Notemos que si $M=0$, sería $f'_s(x) = 0 \ \forall x \in [c,d]$, por lo que $f_s$ sería constante, y la proposición sería cierta para todo $K$ positivo. Es por esto que podemos suponer $M > 0$. Como $s$ era un punto arbitrario de $[a,b]$, queda probado el resultado, donde elegimos $K=M$.
\end{proof}

\textbf{Nota.} En el caso en que $f \in C^1(\mathbb{R}^2,\mathbb{R})$, podemos aplicar directamente el resultado anterior, pues sabemos que en particular $f$ es continua, y derivable respecto de su segunda variable con derivada continua.

\end{ejer}


\end{document}
